\documentclass[11pt,oneside]{book}

\usepackage[T1]{fontenc}
\usepackage[utf8]{inputenc}
\usepackage{lmodern}
\usepackage{microtype}
\usepackage[margin=1.05in,headheight=25pt]{geometry}
\usepackage{amsmath,amssymb,amsthm,mathtools}
% Canonical mathematical notation for active FabiusFunction LaTeX documents.
%
% This file is intentionally a plain TeX input rather than a LaTeX package.
% Every standalone document loads it by an explicit path relative to that
% document's directory; included fragments inherit the definitions from their
% root document.  Keep this file independent of document layout and metadata.
%
% Contract:
%   * one command has one mathematical meaning throughout the active corpus;
%   * names encode semantic roles rather than a document-local choice of letter;
%   * visually similar but differently normalized objects have different print;
%   * every command below is defined normatively in the notation catalogue.
%
% The active corpus excludes docs/archive/.  Historical sources may retain
% their original notation only inside an explicitly labelled quotation.

\makeatletter
\@ifundefined{mathbb}{%
  \PackageError{fabius-notation}{The amssymb package must be loaded first}%
    {Load amsmath and amssymb before inputting fabius-notation.tex.}%
}{}
\@ifundefined{operatorname}{%
  \PackageError{fabius-notation}{The amsmath package must be loaded first}%
    {Load amsmath before inputting fabius-notation.tex.}%
}{}
\makeatother

% Version stamp used by the catalogue and by static audits.
\newcommand{\FabiusNotationVersion}{2026-08-31}

% Number systems.  NaturalNumbers is explicitly the nonnegative convention.
\newcommand{\NaturalNumbers}{\mathbb{N}_{0}}
\newcommand{\PositiveIntegers}{\mathbb{N}_{>0}}
\newcommand{\IntegerNumbers}{\mathbb{Z}}
\newcommand{\RationalNumbers}{\mathbb{Q}}
\newcommand{\RealNumbers}{\mathbb{R}}
\newcommand{\ComplexNumbers}{\mathbb{C}}
\newcommand{\ExtendedRealNumbers}{\overline{\mathbb{R}}}

% The three principal Fabius--Rvachev functions.  Subscripts are deliberate:
% a bare F and a calligraphic F have several unrelated uses in the corpus.
\newcommand{\FabiusBounded}{F_{\mathrm{bd}}}
\newcommand{\FabiusGlobal}{F_{\mathrm{gl}}}
\newcommand{\FabiusQuantile}{Q_{F}}
\newcommand{\FabiusClampedQuantile}{Q_{F,\mathrm{cl}}}
\newcommand{\RvachevUp}{\operatorname{up}}

% Constants, elementary operators, and delimiters.
\newcommand{\EulerE}{\mathrm{e}}
\newcommand{\ImaginaryUnit}{\mathrm{i}}
\newcommand{\Differential}{\mathop{}\!\mathrm{d}}
\newcommand{\AbsoluteValue}[1]{\left\lvert #1\right\rvert}
\newcommand{\Norm}[1]{\left\lVert #1\right\rVert}
\newcommand{\Floor}[1]{\left\lfloor #1\right\rfloor}
\newcommand{\Ceiling}[1]{\left\lceil #1\right\rceil}
\newcommand{\FractionalPart}[1]{\left\{#1\right\}}
\newcommand{\PositivePart}[1]{\left(#1\right)_{+}}
\newcommand{\IndicatorOf}[1]{\mathbf{1}_{#1}}
\newcommand{\SetOf}[1]{\left\{#1\right\}}
\newcommand{\InnerProduct}[1]{\left\langle #1\right\rangle}
\newcommand{\CardinalityOf}[1]{\left\lvert #1\right\rvert}
\newcommand{\DefinitionEquals}{\mathrel{\coloneqq}}
\newcommand{\Sign}{\operatorname{sgn}}
\newcommand{\IdentityOperator}{\operatorname{id}}
\newcommand{\SpanOperator}{\operatorname{span}}
\newcommand{\TraceOperator}{\operatorname{tr}}
\newcommand{\RankOperator}{\operatorname{rank}}
\newcommand{\DiagonalOperator}{\operatorname{diag}}
\newcommand{\ResidueOperator}{\operatorname*{Res}}
\newcommand{\LeastCommonMultiple}{\operatorname{lcm}}
\newcommand{\OrderOperator}{\operatorname{ord}}
\newcommand{\PrincipalLogarithm}{\operatorname{Log}}
\newcommand{\FallingFactorial}[2]{\left(#1\right)^{\underline{#2}}}
\newcommand{\RisingFactorial}[2]{\left(#1\right)^{\overline{#2}}}

% Support, topology, convolution, and metric notation.
\newcommand{\SupportOperator}{\operatorname{supp}}
\newcommand{\SupportOf}[1]{\SupportOperator\!\left(#1\right)}
\newcommand{\TopologicalSupportOf}[1]{\operatorname{tsupp}\!\left(#1\right)}
\newcommand{\Convolution}{\mathbin{*}}
\newcommand{\MetricDistance}{\operatorname{dist}}
\newcommand{\TotalVariationDistance}{d_{\mathrm{TV}}}

% Probability.  Equality and convergence in law are not metric distance.
\newcommand{\Probability}{\mathbb{P}}
\newcommand{\Expectation}{\mathbb{E}}
\newcommand{\Variance}{\operatorname{Var}}
\newcommand{\Covariance}{\operatorname{Cov}}
\newcommand{\LawOperator}{\operatorname{Law}}
\newcommand{\LawOf}[1]{\LawOperator\!\left(#1\right)}
\newcommand{\EqualInLaw}{\mathrel{\overset{\mathrm{law}}{=}}}
\newcommand{\ConvergesInLaw}{\mathrel{\xrightarrow{\mathrm{law}}}}

% Fourier and integral transforms.  The printed labels expose normalization.
% FourierTwoPi uses exp(-2*pi*i*x*xi); FourierAngular uses exp(-i*x*omega).
\newcommand{\FourierTwoPi}[1]{\widehat{#1}_{\,2\pi}}
\newcommand{\FourierAngular}[1]{\widehat{#1}_{\,\mathrm{ang}}}
\newcommand{\LaplaceTransformSymbol}{\mathcal{L}}
\newcommand{\MellinTransformSymbol}{\mathcal{M}}
\newcommand{\LaplaceTransformOf}[1]{\LaplaceTransformSymbol\!\left[#1\right]}
\newcommand{\MellinTransformOf}[1]{\MellinTransformSymbol\!\left[#1\right]}
\newcommand{\MomentGeneratingFunction}[1]{M_{#1}}
\newcommand{\CharacteristicFunction}[1]{\varphi_{#1}}

% The two standard sinc functions must never share one symbol.
\newcommand{\SincRad}{\operatorname{sinc}_{\mathrm{rad}}}
\newcommand{\SincPi}{\operatorname{sinc}_{\pi}}

% Binary arithmetic and Thue--Morse notation.
\newcommand{\BinaryDigitSum}{\operatorname{wt}_{2}}
\newcommand{\ThueMorseSignSymbol}{\varepsilon}
\newcommand{\ThueMorseSign}[1]{\varepsilon_{#1}}
\newcommand{\PadicValuation}[1]{\nu_{#1}}
\newcommand{\TwoAdicValuation}{\nu_{2}}

% q-series notation.  Every base is an explicit argument.
\newcommand{\QPochhammer}[3]{\left(#1;#2\right)_{#3}}
\newcommand{\GaussianBinomial}[3]{%
  \genfrac{[}{]}{0pt}{}{#1}{#2}_{#3}%
}
\newcommand{\QInteger}[2]{\left[#1\right]_{#2}}

% Named special functions.
\newcommand{\Polylogarithm}{\operatorname{Li}}
\newcommand{\LambertW}{W}
\newcommand{\GammaFunction}{\Gamma}
\newcommand{\RiemannZeta}{\zeta}

% Asymptotics.  BigO and LittleO are operator tokens for existing formula
% grammar; prefer the At forms so that the limiting regime is printed.
\newcommand{\BigO}{\mathcal{O}}
\newcommand{\LittleO}{o}
\newcommand{\BigOAt}[2]{\mathcal{O}_{#2}\!\left(#1\right)}
\newcommand{\LittleOAt}[2]{o_{#2}\!\left(#1\right)}

\usepackage{bm}
\usepackage{booktabs,longtable,tabularx,array,multirow}
\usepackage{enumitem}
\usepackage{xcolor}
\usepackage{hyperref}
\usepackage[nameinlink,capitalise,noabbrev]{cleveref}
\usepackage{fancyhdr}
\usepackage{listings}
\usepackage{fancyvrb}
\usepackage{graphicx}
\usepackage{caption}
\usepackage{etoolbox}

\definecolor{linkblue}{RGB}{25,78,135}
\definecolor{deepgreen}{RGB}{22,105,70}
\definecolor{deepred}{RGB}{150,45,45}
\definecolor{softgray}{RGB}{245,246,248}
\hypersetup{
  colorlinks=true,
  linkcolor=linkblue,
  citecolor=deepgreen,
  urlcolor=linkblue,
  pdftitle={Inverse Functions of q-Pochhammer Symbols, Gaussian Coefficients, and Related q-Analogs},
  pdfauthor={Research report for the ProveIt Fabius/q-series program}
}

\pagestyle{fancy}
\fancyhf{}
\fancyhead[L]{\small Inverse functions of $q$-analogs}
\fancyhead[R]{}
\fancyfoot[C]{\thepage}
\renewcommand{\headrulewidth}{0.4pt}

\setcounter{tocdepth}{2}
\setcounter{secnumdepth}{3}
\setlist[itemize]{topsep=4pt,itemsep=2pt,parsep=1pt}
\setlist[enumerate]{topsep=4pt,itemsep=2pt,parsep=1pt}
\allowdisplaybreaks

\newtheorem{theorem}{Theorem}[chapter]
\newtheorem{proposition}[theorem]{Proposition}
\newtheorem{lemma}[theorem]{Lemma}
\newtheorem{corollary}[theorem]{Corollary}
\newtheorem{conjecture}[theorem]{Conjecture}
\newtheorem{problem}[theorem]{Problem}
\theoremstyle{definition}
\newtheorem{definition}[theorem]{Definition}
\newtheorem{example}[theorem]{Example}
\newtheorem{algorithm}[theorem]{Algorithm}
\theoremstyle{remark}
\newtheorem{remark}[theorem]{Remark}

\newcommand{\Disc}{\operatorname{Disc}}
\newcommand{\Cat}{\operatorname{Cat}}
\newcommand{\OO}{\mathcal{O}}
\newcommand{\eps}{\varepsilon}
\newcommand{\statusproved}{\textcolor{deepgreen}{\textbf{proved here}}}
\newcommand{\statusderived}{\textcolor{linkblue}{\textbf{derived here}}}
\newcommand{\statusconj}{\textcolor{deepred}{\textbf{conjectural}}}

\lstdefinestyle{pythonstyle}{
  language=Python,
  basicstyle=\ttfamily\footnotesize,
  keywordstyle=\bfseries,
  commentstyle=\itshape,
  showstringspaces=false,
  breaklines=true,
  frame=single,
  backgroundcolor=\color{softgray},
  columns=fullflexible
}

\title{\Huge Inverse Functions of $q$-Pochhammer Symbols,\\[4pt]
Gaussian Coefficients, and Related $q$-Analogs\\[12pt]
\Large Local Reversion, Puiseux Branches, Singular Limits,\\
Root-of-Unity Scaling, and Global Real Inverses}
\author{Research report for the ProveIt Fabius/$q$-series program}
\date{30 August 2026}

\begin{document}
\hypersetup{pageanchor=false}
\maketitle
\hypersetup{pageanchor=true}
\frontmatter

\chapter*{Abstract}
\addcontentsline{toc}{chapter}{Abstract}
This report develops a systematic inverse-function theory for finite and
infinite $q$-Pochhammer symbols, Gaussian $q$-binomial coefficients,
Jackson's $q$-gamma function, complex-order shifted factorials, and several
related $q$-analogs.  The independent variables being inverted include the
base $q$, the Pochhammer parameter $a$, the complex order $\alpha$, the two
continuous Gaussian parameters, and the arguments of $q$-exponentials,
$q$-numbers, Rogers--Szeg\H{o} polynomials, theta products, and a normalized
$q$-Lambert map.

The guiding point is that ``the inverse'' is normally not a single function.
It is a branched analytic correspondence whose sheets meet above critical
values.  Ordinary Taylor reversion applies at regular points; Puiseux series
appear at multiple zeros and critical points; logarithmic and transseries-like
coordinates are required at the singular boundary $q\to 1$; and algebraic
coordinates are natural at $q\to\pm\infty$ for finite $q$-products.

Among the concrete results are: explicit inverse series for $(a;q)_n$ and
$(a;q)_\infty$ in $a$ and $q$; a stable coefficient law near $q=0$ together
with every low-order exception through fifth order; logarithmic inversion at
$q=1$; a complete Puiseux construction at $q=\infty$; exact classification of
the exceptional inverse branches at $q=-1$; a root-of-unity radial theorem
\[
 \log(a;\zeta e^{-t})_\infty
 =-\frac{\Polylogarithm_2(a^d)}{d^2t}
 +\sum_{r=0}^{d-1}\left(\frac12-\frac rd\right)
      \PrincipalLogarithm(1-a\zeta^r)+\OO(t),
\]
for primitive $d$-th roots $\zeta$; the resulting $a\asymp t^{1/d}$ inverse
law; global positive-real inverses of Gaussian polynomials; expansions at
$q=0,1,\pm\infty$ and roots of unity; exact value and derivative formulas at
$q=-1$; inverse theory for continuous Gaussian parameters via $q$-digamma
functions; and $q\to0,1,\infty$ inversion for $\Gamma_q$.

A companion Python program performs high-precision numerical and exact
symbolic checks.  Conjectures are explicitly separated from proved or formally
derived statements.  The final chapters formulate research programs on
monodromy and generic Galois groups, resurgent corrections at the unit circle,
certified continuation, coefficient positivity, and formal verification.


\chapter*{Reader's guide and status conventions}
\addcontentsline{toc}{chapter}{Reader's guide and status conventions}
The target repository subtree contains two substantial source documents: a
frontier volume on exponents and $q$-series, and a proof-oriented monograph on
$q$-Pochhammer symbols and Gaussian coefficients
\cite{ProveItFrontiers,ProveItMonograph}.  They establish a broad \emph{direct}
theory: finite and infinite products, logarithmic derivatives, $q$-gamma
representations, $q$-Newton interpolation, root-of-unity algebra, and
$q\to1$ asymptotics.  They do not organize these facts into a systematic
theory of functional inversion.  The present report supplies that missing
layer and deliberately refers back to the direct theory rather than reproducing
its many identity proofs.

Statements marked \statusproved\ have proofs in this report.  Formulae marked
\statusderived\ follow by explicit formal or asymptotic calculation; their
analytic domains are stated where needed.  Conjectures and open problems are
marked \statusconj.  Numerical experiments are consistency checks, never
substitutes for proof.

The main recurring notation is
\[
 (a;q)_n=\prod_{j=0}^{n-1}(1-aq^j),\qquad
 (a;q)_\infty=\prod_{j=0}^{\infty}(1-aq^j),
\]
\[
 \GaussianBinomial{n}{k}{q}=\frac{(q;q)_n}{(q;q)_k(q;q)_{n-k}},\qquad
 [x]_q=\frac{1-q^x}{1-q}.
\]
For singular limits at $q=1$ we normally write $q=e^{-t}$ or $q=e^t$,
according to whether the function is naturally defined inside the unit disk.
Branches of $\PrincipalLogarithm$, fractional powers, and polylogarithms are always part of
the data of a complex inverse branch.

\tableofcontents

\mainmatter

\part{General architecture of inverse $q$-theory}

\chapter{Source audit, scope, and the inverse correspondence}

\section{What has to be inverted}
A $q$-analog usually depends on several parameters, so every direct identity
creates several different inverse problems.  For the Pochhammer symbol alone
one may ask for
\[
 a=A_n(y;q),\qquad q=Q_n(y;a),\qquad
 \alpha=\mathcal A(y;a,q)
\]
from $y=(a;q)_n$ or $y=(a;q)_\alpha$.  The first equation is algebraic of
degree $n$ when $n<\infty$; the second is algebraic of degree
$n(n-1)/2$; and the third is transcendental but monotone on a useful real
domain.  No single notation can hide the different branch geometries.

For a Gaussian polynomial, $y=\GaussianBinomial{n}{k}{q}$ is an algebraic equation of
degree $D=k(n-k)$ in $q$.  If $n,k$ are replaced by continuous parameters
$\alpha,\beta$ through $\Gamma_q$, inversion in $\alpha$ is regular and
one-sided on the positive real domain, while inversion in $\beta$ has a
symmetry-forced square-root branch at $\beta=\alpha/2$.

\section{Inverse functions as Riemann surfaces}
Let $F(x;\lambda)$ be holomorphic in $x$ and in auxiliary parameters
$\lambda$.  The natural inverse object is
\[
 \mathcal X_F=\{(x,y,\lambda):F(x;\lambda)=y\}
 \longrightarrow \{(y,\lambda)\}.
\]
Regular sheets are locally graphs $x=X(y;\lambda)$.  They meet where
$F_x=0$ or where the direct function itself becomes singular.  In polynomial
problems the finite branch locus is controlled by
\[
 \Disc_x(F(x;\lambda)-y)
 =\ResidueOperator_x(F(x;\lambda)-y,F_x(x;\lambda)).
\]
This viewpoint prevents three common errors: treating a multivalued inverse as
single-valued, using Taylor series at a critical point, and confusing a zero of
$F$ with a branch point of $F^{-1}$.  A simple zero gives a regular inverse;
a multiple zero gives a Puiseux branch.

\section{A map of the principal regimes}
\begin{longtable}{>{\raggedright\arraybackslash}p{0.21\textwidth}
                  >{\raggedright\arraybackslash}p{0.20\textwidth}
                  >{\raggedright\arraybackslash}p{0.24\textwidth}
                  >{\raggedright\arraybackslash}p{0.245\textwidth}}
\toprule
Direct object and variable & Base point & Natural inverse coordinate & Typical geometry\\
\midrule
\endfirsthead
\toprule
Direct object and variable & Base point & Natural inverse coordinate & Typical geometry\\
\midrule
\endhead
$(a;q)_n$ in $a$ & $a=0$ & $1-y$ & Taylor, unless $[n]_q=0$\\
$(a;q)_n$ in $a$ & $a=q^{-r}$ & $y^{1/\mu}$ & simple or multiple-zero Puiseux\\
$(a;q)_n$ in $q$ & $q=0$ & $\{y-(1-a)\}/(a^2-a)$ & stable Taylor reversion\\
$(a;q)_n$ in $q$ & $q=1$ & $\log\{y/(1-a)^n\}$ & logarithmic Taylor reversion\\
$(a;q)_n$ in $q$ & $q=\pm\infty$ & $y^{-1/N}$ & $N$ Puiseux sheets\\
$(a;q)_\infty$ in $q$ & $q\to\zeta$, $|q|<1$ & $1/(-\log y)$ & dilogarithmic singular inversion\\
$(a;q)_\infty$ in $a$ & $q\to\zeta$ & $t^{1/d}$ & root-of-unity Puiseux scaling\\
$\GaussianBinomial{n}{k}{q}$ in $q$ & $q=0$ & $y-1$ & partition-coefficient reversion\\
$\GaussianBinomial{n}{k}{q}$ in $q$ & $q=1$ & $\log\{y/\binom nk\}$ & Bernoulli-cumulant reversion\\
$\GaussianBinomial{n}{k}{q}$ in $q$ & $q=\pm\infty$ & $y^{-1/D}$ & reciprocal Puiseux sheets\\
continuous $q$-binomial in $\beta$ & $\beta=\alpha/2$ & $\sqrt{-\log(y/y_{\max})}$ & symmetry-forced double branch\\
$\Gamma_q(x)$ in $x$ & minimum & $\sqrt{y-y_{\min}}$ & two real sheets\\
\bottomrule
\end{longtable}

\chapter{A reusable calculus of reversion and ramification}

\section{Regular local inversion}
\begin{theorem}[Holomorphic inverse and coefficient reversion]\label{thm:regular-inverse}
Let
\[
 u=f(x)=c_1x+\frac{c_2}{2!}x^2+\frac{c_3}{3!}x^3
       +\frac{c_4}{4!}x^4+\cdots,
 \qquad c_1\ne0.
\]
Then the branch satisfying $x(0)=0$ is holomorphic and begins
\begin{align}
 x={}&\frac{u}{c_1}-\frac{c_2}{2c_1^3}u^2
 +\frac{3c_2^2-c_1c_3}{6c_1^5}u^3 \notag\\
 &+\frac{-c_1^2c_4+10c_1c_2c_3-15c_2^3}
 {24c_1^7}u^4+\OO(u^5).\label{eq:universal-reversion}
\end{align}
More generally, if $f(x)=x/\phi(x)$ with $\phi(0)\ne0$, then
\[
 [u^m]x(u)=\frac1m[z^{m-1}]\phi(z)^m.
\]
\end{theorem}
\begin{proof}
The analytic assertion is the inverse-function theorem.  The coefficient
formula is Lagrange--B\"urmann inversion.  Substituting a formal ansatz
$x=b_1u+b_2u^2+\cdots$ and equating coefficients gives
\eqref{eq:universal-reversion}; the calculation takes place in the complete
local ring $\ComplexNumbers[[u]]$ and therefore also proves formal uniqueness.
\end{proof}

\begin{remark}
For $q$-products it is usually better to revert a logarithm.  If $F(x_0)\ne0$,
set $u=\PrincipalLogarithm(F(x)/F(x_0))$.  Logarithmic derivatives turn products into Lambert
sums, and the $c_j$ in \eqref{eq:universal-reversion} become cumulants rather
than raw derivatives.
\end{remark}

\section{Critical points and Puiseux inversion}
\begin{theorem}[Local Puiseux law]\label{thm:puiseux}
Suppose
\[
 F(x)-y_0=c_m(x-x_0)^m+c_{m+1}(x-x_0)^{m+1}+\cdots,
 \qquad c_m\ne0,
\]
with $m\ge2$.  Then the inverse correspondence has $m$ local sheets and each
has a convergent Puiseux expansion
\[
 x-x_0=\omega\left(\frac{y-y_0}{c_m}\right)^{1/m}
 +\OO\bigl((y-y_0)^{2/m}\bigr),\qquad \omega^m=1.
\]
\end{theorem}
\begin{proof}
Write $F(x)-y_0=(x-x_0)^mG(x)$ with $G(x_0)=c_m\ne0$.  Choose a local
holomorphic $m$-th root $G(x)^{1/m}$ and apply the regular inverse theorem to
$(x-x_0)G(x)^{1/m}$.
\end{proof}

The theorem covers multiple Pochhammer zeros, symmetry-forced extrema, and
root-of-unity coalescence.  Newton polygons provide the analogous construction
when two small parameters compete, as in $a\to0$ and $q\to\zeta$.

\section{Implicit derivatives with respect to every parameter}
If $F(x,p)=y$ and $F_x\ne0$, the inverse branch $x=X(y,p)$ satisfies
\begin{align}
 X_y&=\frac1{F_x},&
 X_p&=-\frac{F_p}{F_x},\label{eq:implicit-first}\\
 X_{yy}&=-\frac{F_{xx}}{F_x^3},&
 X_{pp}&=-\frac{F_{pp}+2F_{xp}X_p+F_{xx}X_p^2}{F_x}.
 \label{eq:implicit-second}
\end{align}
Mixed higher derivatives are efficiently encoded by multivariate Bell
polynomials.  Equations \eqref{eq:implicit-first}--\eqref{eq:implicit-second}
are the basic engine for expansions in $q$, $a$, order, or any auxiliary
parameter without re-expanding the whole product.

\section{Singular asymptotic inversion}
Suppose $Y\to\infty$ and
\[
 Y=\frac{A}{t}-B+Ct-Et^3+\OO(t^5),\qquad A\ne0.
\]
With $\Lambda=Y+B$, direct reversion gives
\begin{equation}\label{eq:singular-master}
 t=\frac{A}{\Lambda}
 +\frac{A^2C}{\Lambda^3}
 +\frac{2A^3C^2-A^4E}{\Lambda^5}
 +\OO(\Lambda^{-7}).
\end{equation}
This simple template will invert the McIntosh expansion of an infinite
$q$-Pochhammer product at $q=1$, and its root-of-unity analogues.

\section{Conditioning and rigorous residual bounds}
For a real monotone branch, if $|F'(x)|\ge m>0$ between an approximation
$\widehat x$ and the exact root $x_*$, then
\[
 |\widehat x-x_*|\le \frac{|F(\widehat x)-y|}{m}.
\]
Thus every truncated inverse series can be turned into a certified enclosure
once a derivative lower bound is available.  In the complex plane, a Rouch\'e
argument on a disk or an interval-Newton/Krawczyk operator gives the analogous
certificate.  The inverse condition number is
\[
 \kappa_{\mathrm{abs}}=\frac1{|F_x|},\qquad
 \kappa_{\mathrm{rel}}=\left|\frac{y}{xF_x}\right|,
\]
and diverges precisely as a regular sheet approaches the critical locus.

\part{Finite $q$-Pochhammer symbols}

\chapter{Inversion in the Pochhammer parameter $a$}

Throughout this chapter
\[
 P_n(a,q):=(a;q)_n=\prod_{j=0}^{n-1}(1-aq^j),\qquad n\ge1.
\]
As a polynomial in $a$, it has degree $n$ and zeros $q^{-r}$,
$0\le r<n$, when $q\ne0$.  Hence a generic level equation
$P_n(a,q)=y$ has $n$ complex sheets.

\section{A global real inverse}
\begin{theorem}[Global $a$-inverse for $0<q<1$]\label{thm:global-a-finite}
Fix $0<q<1$ and $n\ge1$.  The map
\[
 a\longmapsto P_n(a,q)
\]
is positive and strictly decreasing from $+\infty$ to $0$ on
$(-\infty,1)$.  Consequently, for every $y>0$ there is a unique real solution
$a<1$ of $P_n(a,q)=y$.
\end{theorem}
\begin{proof}
For $a<1$ every factor is positive and
\[
 \partial_a\log P_n(a,q)
 =-\sum_{j=0}^{n-1}\frac{q^j}{1-aq^j}<0.
\]
The limit at $a\to1^-$ is $0$.  As $a\to-\infty$, the leading term
$(-a)^nq^{n(n-1)/2}$ tends to $+\infty$.
\end{proof}

This theorem gives a canonical real sheet.  The remaining $n-1$ sheets are
essential for complex continuation and for the $q\to\infty$ analysis below.

\section{Reversion at $a=0$}
The finite $q$-binomial theorem gives
\begin{equation}\label{eq:finite-a-expansion}
 P_n(a,q)=\sum_{r=0}^n(-a)^r E_r(n,q),
 \qquad E_r(n,q)=q^{r(r-1)/2}\GaussianBinomial{n}{r}{q}.
\end{equation}
Thus $E_1=[n]_q$.  Put $z=1-y$.

\begin{proposition}[Inverse at the unit level]\label{prop:a-inverse-zero}
If $E_1(n,q)\ne0$, the branch with $a(1)=0$ is
\begin{align}
 a={}&\frac{z}{E_1}+\frac{E_2}{E_1^3}z^2
 +\left(\frac{2E_2^2}{E_1^5}-\frac{E_3}{E_1^4}\right)z^3
 \notag\\
 &+\left(\frac{5E_2^3}{E_1^7}
 -\frac{5E_2E_3}{E_1^6}+\frac{E_4}{E_1^5}\right)z^4
 +\OO(z^5).\label{eq:a-inverse-zero}
\end{align}
\end{proposition}
\begin{proof}
Equation \eqref{eq:finite-a-expansion} says
$z=E_1a-E_2a^2+E_3a^3-E_4a^4+\cdots$.  Apply
\cref{thm:regular-inverse}, or equate coefficients directly.
\end{proof}

At $q=1$, $E_r=\binom nr$ and \eqref{eq:a-inverse-zero} sums to the exact
branch $a=1-y^{1/n}$.  At a nontrivial root of unity, the hypothesis can fail
in a structured way.

\begin{theorem}[Root-of-unity ramification at $a=0$]\label{thm:a-zero-root-unity}
Let $\zeta$ be a primitive $d$-th root of unity and write $n=md+r$ with
$0\le r<d$.
\begin{enumerate}
\item If $r>0$, then $[n]_\zeta=[r]_\zeta\ne0$ and the inverse at $a=0$
      is regular.
\item If $r=0$, then
\[
 P_n(a,\zeta)=(1-a^d)^m=1-ma^d+\OO(a^{2d}),
\]
so the inverse has $d$ Puiseux sheets
\[
 a=\omega\left(\frac{1-y}{m}\right)^{1/d}
 +\OO\bigl((1-y)^{2/d}\bigr),\qquad \omega^d=1.
\]
\end{enumerate}
\end{theorem}
\begin{proof}
Group the factors into complete residue blocks modulo $d$ and use
$\prod_{j=0}^{d-1}(1-a\zeta^j)=1-a^d$.  If a partial block remains, its
linear coefficient is $-[r]_\zeta\ne0$.
\end{proof}

\section{Inversion at the Pochhammer zeros}
Let $a_r=q^{-r}$, $0\le r<n$, and assume $q\ne0$ and that the zeros are
distinct.  Define
\[
 R_{n,r}(q)=\prod_{\substack{0\le j<n\\j\ne r}}(1-q^{j-r})
 =(-1)^rq^{-r(r+1)/2}(q;q)_r(q;q)_{n-1-r}.
\]
Then
\begin{equation}\label{eq:poch-a-derivative-root}
 \partial_aP_n(a_r,q)
 =-q^rR_{n,r}(q)
 =(-1)^{r+1}q^{-r(r-1)/2}(q;q)_r(q;q)_{n-1-r}.
\end{equation}

\begin{proposition}[Local inverse at a simple zero]
Under the preceding hypotheses, the sheet through $a_r$ is
\[
 a=a_r+\frac{y}{P_a(a_r,q)}
 -\frac{P_{aa}(a_r,q)}{2P_a(a_r,q)^3}y^2+\OO(y^3).
\]
\end{proposition}

At a root of unity several zeros may coalesce.

\begin{proposition}[Multiplicity at a root of unity]\label{prop:root-unity-a-mult}
Let $\zeta$ be primitive of order $d$, let $n=md+r$, $0\le r<d$, and let
$a_0=\zeta^{-s}$, $0\le s<d$.  As a function of $a$, the zero of
$P_n(a,\zeta)$ at $a_0$ has multiplicity
\[
 \mu_s=m+\mathbf 1_{\{0\le s<r\}}.
\]
Consequently, the inverse at $y=0$ has $\mu_s$ Puiseux sheets.
\end{proposition}
\begin{proof}
A factor vanishes exactly when its index is congruent to $s$ modulo $d$.
Among $0,\ldots,n-1$, that residue occurs $m$ times and once more precisely
when it belongs to the partial block $0,\ldots,r-1$.
\end{proof}

\section{Exact inverses at $q=0,1,-1$}
Three bases collapse the $a$-equation to elementary algebra.

At $q=0$,
\[
 P_n(a,0)=1-a,
 \qquad a=1-y.
\]
At $q=1$,
\[
 P_n(a,1)=(1-a)^n,
 \qquad a=1-\omega y^{1/n},\quad \omega^n=1.
\]
At $q=-1$, put $u=\lceil n/2\rceil$, $v=\lfloor n/2\rfloor$.  Then
\begin{equation}\label{eq:finite-at-minus-one}
 P_n(a,-1)=(1-a)^u(1+a)^v.
\end{equation}
If $n=2m$, this becomes $(1-a^2)^m$, and the inverse is exactly
\[
 a=\pm\sqrt{1-y^{1/m}},
\]
with all choices of the $m$-th root included in the complex correspondence.
The double branch at $a=0,y=1$ is forced by evenness.  If $n=2m+1$, then
$P_n=(1-a)(1-a^2)^m$ and the sheet at $a=0$ is regular because
$P_a(0,-1)=-1$.

\section{The $a$-inverse for $q$ near $0$ and $1$}
Fix $y$ and regard $a=a(q;y)$.  For the branch satisfying $a(0)=1-y$,
formal implicit reversion gives, for $n\ge5$ through the displayed order,
\begin{align}
 a={}&1-y+y(y-1)q-2y(y-1)^2q^2 \notag\\
 &+y(y-1)^2(5y-3)q^3
 -y(y-1)^2(2y-1)(7y-5)q^4+\OO(q^5).
 \label{eq:a-inverse-q-zero}
\end{align}
The coefficient of $q^M$ stabilizes once $n>M$ because the omitted factors
begin at $q^n$.

For $q=e^t$ near $1$, choose an $n$-th root $r=y^{1/n}$ and the branch
$a(0)=1-r$.  Implicit differentiation yields
\begin{align}
 a(t)={}&1-r-\frac{(n-1)(1-r)}2t\notag\\
 &+\frac{(n-1)(1-r)(3nr-n-3r-1)}{24r}t^2
 +\OO(t^3).\label{eq:a-near-q-one-finite}
\end{align}
The different choices of $r$ parametrize the $n$ complex sheets.

\section{All $a$-branches as $q\to\infty$}
The polynomial equation $P_n(a,q)=y$ has $n$ asymptotic $a$-branches.  This
is more informative than following only the global real branch.

\begin{theorem}[$a$-branch atlas at large base]\label{thm:a-branches-infty}
Let $y\ne0$ be fixed and $q\to\infty$ in a sector.  For
$0\le r\le n-2$, put
\[
 M_r=\frac{(n-1-r)(n-r)}2.
\]
There is a branch near the moving zero $q^{-r}$ with
\begin{equation}\label{eq:a-large-q-moving}
 a_r(q)=q^{-r}\left(1+(-1)^{n-r}yq^{-M_r}
 +\OO(q^{-M_r-1})+\OO(q^{-2M_r})\right).
\end{equation}
The remaining branch has the full expansion
\begin{align}
 a_{n-1}(q)=q^{-(n-1)}\bigl(&1-y+y(y-1)q^{-1}
 -2y(y-1)^2q^{-2}\notag\\
 &+y(y-1)^2(5y-3)q^{-3}+\OO(q^{-4})\bigr).
 \label{eq:a-large-q-last}
\end{align}
The same formal Laurent series describe $q\to-\infty$, with the signs already
encoded by powers of $q^{-1}$.
\end{theorem}
\begin{proof}
For $r\le n-2$, set $a=q^{-r}(1+\eps)$.  The vanishing factor is
$1-aq^r=-\eps$, factors with $j<r$ tend to $1$, and those with $j>r$ have
product $(-1)^{n-1-r}q^{M_r}(1+o(1))$.  Solving the leading equation gives
\eqref{eq:a-large-q-moving}.  For the final branch write
$a=q^{-(n-1)}c$ and $s=q^{-1}$.  Then
\[
 P_n(a,q)=(c;s)_n.
\]
The branch $c(0)=1-y$ is exactly the small-$s$ inverse
\eqref{eq:a-inverse-q-zero}, proving \eqref{eq:a-large-q-last}.
\end{proof}

\chapter{Inversion in the base $q$}

Let
\[
 N=\frac{n(n-1)}2.
\]
For fixed $a$, $P_n(a,q)$ is a polynomial of degree $N$ in $q$, except in
degenerate cases.  Hence a generic base inverse has $N$ complex branches.

\section{A canonical real branch on $0\le q\le1$}
\begin{theorem}[Global finite base inverse]\label{thm:global-q-finite}
Let $n\ge2$.
\begin{enumerate}
\item If $0<a<1$, then $q\mapsto P_n(a,q)$ is strictly decreasing on
$[0,1]$ from $1-a$ to $(1-a)^n$.
\item If $a<0$, it is strictly increasing on $[0,1]$ from $1-a$ to
$(1-a)^n$.
\end{enumerate}
Thus every level strictly between the endpoint values has one and only one
real inverse $q\in(0,1)$.
\end{theorem}
\begin{proof}
For $0<q<1$,
\[
 \partial_q\log P_n(a,q)
 =-a\sum_{j=1}^{n-1}\frac{j q^{j-1}}{1-aq^j}.
\]
All denominators are positive.  The sign is therefore $-\operatorname{sgn}a$.
\end{proof}

The cases $a=0$, $a=1$, $n=0$, and $n=1$ are constant in $q$ and have no
local base inverse.

\section{The stable inverse series at $q=0$}
For $a\notin\{0,1\}$ define the normalized level coordinate
\begin{equation}\label{eq:x-normalization-qzero}
 x=\frac{y-(1-a)}{a^2-a}.
\end{equation}
Then $x=q+\OO(q^2)$.

\begin{theorem}[Stable small-base reversion]\label{thm:stable-qzero}
For the infinite formal product, and for the finite product whenever $n>5$,
\begin{align}
 q={}&x-x^2+(a+1)x^3-(4a+1)x^4\notag\\
 &+(3a^2+11a+1)x^5+\OO(x^6).
 \label{eq:qzero-stable}
\end{align}
More generally, the coefficient through $x^M$ agrees with the infinite-product
coefficient whenever $n>M$.
\end{theorem}
\begin{proof}
The quotient of $(a;q)_\infty$ by $(a;q)_n$ is
$(aq^n;q)_\infty=1+\OO(q^n)$.  Therefore the direct $q$-series agree through
$q^{n-1}$, and formal inversion preserves agreement through the same order.
Direct reversion gives \eqref{eq:qzero-stable}.
\end{proof}

The low-degree exceptions are important enough to record:
\begin{align*}
 n=2:\quad q={}&x,\\
 n=3:\quad q={}&x-x^2+(a+2)x^3-5(a+1)x^4
 +(3a^2+21a+14)x^5+\OO(x^6),\\
 n=4:\quad q={}&x-x^2+(a+1)x^3-4ax^4
 +(3a^2+10a-4)x^5+\OO(x^6),\\
 n=5:\quad q={}&x-x^2+(a+1)x^3-(4a+1)x^4
 +(3a^2+11a+2)x^5+\OO(x^6).
\end{align*}
These identities were also checked exactly in the companion SymPy script.

\section{Logarithmic inversion at $q=1$}
Set $q=e^t$ and
\[
 u=\PrincipalLogarithm\frac{P_n(a,e^t)}{(1-a)^n}.
\]
Let $S_m(n)=\sum_{j=0}^{n-1}j^m$.  Expanding the logarithm of every factor
gives the exact jet
\begin{equation}\label{eq:finite-cumulants-qone}
 u=\sum_{m\ge1}\frac{c_m}{m!}t^m,
 \qquad c_m=-S_m(n)\Polylogarithm_{1-m}(a).
\end{equation}
Thus the inverse $t=t(u)$ is obtained from
\eqref{eq:universal-reversion}.  Explicitly,
\begin{align}
 t={}&\frac{u}{c_1}-\frac{c_2}{2c_1^3}u^2
 +\frac{3c_2^2-c_1c_3}{6c_1^5}u^3+\OO(u^4),
 \label{eq:finite-qone-inverse-t}\\
 q={}&\exp(t).
\end{align}
The first direct derivative is
\[
 \partial_qP_n(a,1)=-aS_1(n)(1-a)^{n-1},
\]
so, in the raw level coordinate,
\begin{equation}\label{eq:finite-qone-linear}
 q-1=-\frac{y-(1-a)^n}{aS_1(n)(1-a)^{n-1}}
 +\OO\bigl((y-(1-a)^n)^2\bigr).
\end{equation}
The inverse is regular unless $a=0$, $a=1$, or $n\le1$.

\section{The complete Puiseux construction at $q=\infty$}
Factor
\begin{equation}\label{eq:finite-q-infty-factor}
 P_n(a,q)=C(a)q^NH(q^{-1}),
\quad C(a)=(1-a)(-a)^{n-1},
\quad H(s)=\prod_{j=1}^{n-1}(1-a^{-1}s^j).
\end{equation}
Assume $a\notin\{0,1\}$ and $y\to\infty$.  Choose an $N$-th root
$x=(C/y)^{1/N}$ and solve
\begin{equation}\label{eq:s-lagrange-infty}
 s=xH(s)^{1/N},\qquad q=s^{-1}.
\end{equation}
Lagrange inversion gives the all-orders coefficient formula
\begin{equation}\label{eq:s-coeff-infty}
 [x^m]s(x)=\frac1m[u^{m-1}]H(u)^{m/N}.
\end{equation}
Every choice of $x$ gives one sheet.  For $n\ge3$ the first terms are
\begin{equation}\label{eq:q-infty-finite-first}
 q=x^{-1}+\frac1{aN}
 +\left(\frac1{aN}+\frac{N-1}{2N^2a^2}\right)x+\OO(x^2).
\end{equation}
For $n=2$, this reduces to the exact linear inverse
$q=x^{-1}+a^{-1}$.  The negative-real asymptotic $q\to-\infty$ is not a
separate formula: it is the sheet for which the leading root $x^{-1}$ is
negative real.  The parity of $N$ and the sign of $C(a)$ determine which real
levels reach that sheet.

\section{Zeros in the $q$-plane}
For fixed $a\ne0$, every equation
\[
 q^r=a^{-1},\qquad 1\le r<n,
\]
produces zeros of $P_n$ as a function of $q$.  If $q_0$ makes only the
$r$-th factor vanish, then
\begin{equation}\label{eq:qzero-derivative-general}
 P_q(a,q_0)=-arq_0^{r-1}
 \prod_{\substack{0\le j<n\\j\ne r}}(1-aq_0^j),
\end{equation}
and the inverse at $y=0$ is regular.  If several indices satisfy
$aq_0^j=1$, their number is the zero multiplicity and
\cref{thm:puiseux} applies.  Such collisions occur precisely on arithmetic
strata involving roots of unity.

\section{The exceptional geometry at $q=-1$}
For $a\ne\pm1$, the logarithmic base derivative at $q=-1$ is
\begin{equation}\label{eq:logder-minusone}
 \frac{P_q(a,-1)}{P_n(a,-1)}
 =a\left(\frac{S_{\mathrm{even}}}{1-a}
          -\frac{S_{\mathrm{odd}}}{1+a}\right),
\end{equation}
where the sums run over $0\le j<n$.  Evaluation gives the following complete
classification apart from constant degeneracies.

\begin{theorem}[Critical Pochhammer parameters at $q=-1$]\label{thm:critical-minus-one}
\begin{enumerate}
\item If $n=2m$, then
\[
 \frac{P_q}{P}=\frac{am((2m-1)a-1)}{1-a^2}.
\]
For $m\ge2$, the nontrivial critical parameter is
$a_c=(2m-1)^{-1}$ and
\[
 \left.\frac{P_{qq}}P\right|_{(a_c,-1)}
 =-\frac{(m+1)(2m-1)^2}{12m(m-1)}\ne0.
\]
The inverse therefore has a square-root branch at
$y_c=P_n(a_c,-1)$.
\item If $n=2m+1$, then
\[
 \frac{P_q}{P}=\frac{am(1+(2m+1)a)}{1-a^2}.
\]
For $m\ge2$, the critical parameter is $a_c=-(2m+1)^{-1}$ and
\[
 \left.\frac{P_{qq}}P\right|_{(a_c,-1)}
 =-\frac{(m-1)(2m+1)^2}{12m(m+1)}\ne0.
\]
Again the inverse is square-root ramified.
\item The exceptional case $n=3$, $a=-1/3$ is cubic rather than quadratic:
\[
 P_3(-1/3,q)=\frac{32}{27}+\frac4{27}(q+1)^3.
\]
Hence
\[
 q=-1+\left(\frac{27}{4}\left(y-\frac{32}{27}\right)\right)^{1/3}
\]
with three local branches.
\end{enumerate}
\end{theorem}
\begin{proof}
Equation \eqref{eq:logder-minusone} follows by differentiating the product.
The even and odd arithmetic sums give the displayed first derivatives.
A second differentiation and elementary power sums give the two ratios for
$P_{qq}/P$.  The cubic identity in the last case is an exact expansion.
\end{proof}

\section{General roots of unity}
Let $q=\zeta$ be primitive of order $d$.  The direct product value is
\[
 P_n(a,\zeta)=(1-a^d)^m(a;\zeta)_r,
 \qquad n=md+r,\quad 0\le r<d.
\]
The derivative criterion
\[
 -a\sum_{j=1}^{n-1}\frac{j\zeta^{j-1}}{1-a\zeta^j}\ne0
\]
ensures an ordinary local $q$-inverse.  Its vanishing defines an explicit
algebraic critical locus in $a$.  When a factor also vanishes, one first
extracts all zero factors and applies the Newton polygon to the residual
two-variable germ.  This produces a finite, algorithmic classification of
root-of-unity branches for every fixed $n$.

\part{Infinite products and complex order}

\chapter{The infinite $q$-Pochhammer inverse}

Assume initially $|q|<1$.  The product
\[
 P_\infty(a,q):=(a;q)_\infty
\]
is entire in $a$, while its dependence on $q$ is naturally tied to the unit
disk.  The Lambert-series identity
\begin{equation}\label{eq:lambert-log}
 \PrincipalLogarithm(a;q)_\infty
 =-\sum_{m\ge1}\frac{a^m}{m(1-q^m)}
\end{equation}
holds absolutely for $|a|<1$, $|q|<1$ and extends by analytic continuation in
$a$ away from chosen logarithmic cuts.  It is the central tool for inverse
asymptotics.

\section{Global real inverses in $a$ and $q$}
\begin{theorem}[Global real sheets for the infinite product]\label{thm:global-infinite}
Fix $0<q<1$.  The function $a\mapsto(a;q)_\infty$ is positive and strictly
decreasing from $+\infty$ to $0$ on $(-\infty,1)$, so it has a unique real
$a$-inverse for every $y>0$.

Fix instead $a\in(-\infty,1)\setminus\{0\}$.  On $0\le q<1$ the function is
strictly decreasing when $0<a<1$ and strictly increasing when $a<0$.  Its
ranges are
\[
 0<a<1:\quad (0,1-a],
 \qquad
 a<0:\quad [1-a,\infty).
\]
Hence it has a unique real base inverse on the corresponding range.
\end{theorem}
\begin{proof}
The $a$-derivative is
\[
 \partial_a\PrincipalLogarithm(a;q)_\infty
 =-\sum_{j\ge0}\frac{q^j}{1-aq^j}<0
\]
for $a<1$.  The large negative-$a$ limit follows, for example, by retaining an
increasing number of factors with $|a|q^j\ge1$.  The base derivative has the
same sign calculation as in \cref{thm:global-q-finite}.  Finally,
\eqref{eq:lambert-log} and the $q\to1$ asymptotic below show convergence to
$0$ for $0<a<1$ and divergence to $+\infty$ for $a<0$.
\end{proof}

\section{The $a$-inverse at $a=0$}
Euler's expansion gives
\[
 (a;q)_\infty=\sum_{r\ge0}
 \frac{(-1)^rq^{r(r-1)/2}}{(q;q)_r}a^r.
\]
Substitution into \eqref{eq:a-inverse-zero} yields an explicit branch.  With
$z=1-y$,
\begin{align}
 a={}&(1-q)z+\frac{q(1-q)}{1+q}z^2\notag\\
 &+\frac{q^2(1-q)(q^2+q+2)}{(1+q)^2(1+q+q^2)}z^3
 +\OO(z^4).\label{eq:ainf-near-zero}
\end{align}
This expansion is valid away from roots of the displayed cyclotomic
denominators, and elsewhere after the appropriate Puiseux rescaling.

\section{The $a$-inverse at the first zero $a=1$}
Let
\[
 C(q)=(q;q)_\infty,
 \qquad S_r(q)=\sum_{j\ge1}\left(\frac{q^j}{1-q^j}\right)^r,
 \qquad w=\frac{y}{C(q)}.
\]
Writing $z=1-a$ gives
\[
 (a;q)_\infty=C(q)z
 \exp\left(S_1z-\frac{S_2}{2}z^2+\frac{S_3}{3}z^3-\cdots\right).
\]
Therefore
\begin{equation}\label{eq:ainf-near-one}
 1-a=w-S_1w^2+\frac{3S_1^2+S_2}{2}w^3+\OO(w^4).
\end{equation}
At the general zero $a_r=q^{-r}$, $r\ge0$, one has
\begin{equation}\label{eq:infinite-a-derivative-root}
 \partial_a(a;q)_\infty\big|_{a=q^{-r}}
 =(-1)^{r+1}q^{-r(r-1)/2}(q;q)_r(q;q)_\infty,
\end{equation}
so every zero is simple when $q$ is not a root of unity and the local inverse
is ordinary Taylor.

\section{The base inverse at $q=0$}
The stable finite formula \eqref{eq:qzero-stable} is exactly the infinite
formula:
\[
 x=\frac{y-(1-a)}{a^2-a},
\]
\begin{equation}\label{eq:qinftyprod-qzero}
 q=x-x^2+(a+1)x^3-(4a+1)x^4
 +(3a^2+11a+1)x^5+\OO(x^6).
\end{equation}
Conversely, fixing $y$ and inverting in $a$ gives the all-product version of
\eqref{eq:a-inverse-q-zero}:
\begin{align}
 a={}&1-y+y(y-1)q-2y(y-1)^2q^2\notag\\
 &+y(y-1)^2(5y-3)q^3
 -y(y-1)^2(2y-1)(7y-5)q^4+\OO(q^5).
 \label{eq:ainf-qzero-fixed-y}
\end{align}

\section{Singular base inversion at $q\to1^-$}
Put $q=e^{-t}$, $t\downarrow0$, and fix $a$ in a compact subset of a domain
where the chosen branches are analytic.  The McIntosh expansion
\cite{McIntosh1999,DLMFq} is
\begin{align}
 \PrincipalLogarithm(a;e^{-t})_\infty
={}&-\frac{\Polylogarithm_2(a)}{t}+\frac12\PrincipalLogarithm(1-a)\notag\\
 &-\sum_{r=1}^{M-1}\frac{B_{2r}}{(2r)!}
 t^{2r-1}\Polylogarithm_{2-2r}(a)+\OO(t^{2M-1}).
 \label{eq:mcintosh}
\end{align}
For $0<a<1$, let
\[
 Y=-\log y,\quad
 A=\Polylogarithm_2(a),\quad B=\frac12\log(1-a),
\]
\[
 C=\frac{\Polylogarithm_0(a)}{12}=\frac{a}{12(1-a)},
 \qquad E=\frac{\Polylogarithm_{-2}(a)}{720}
 =\frac{a(1+a)}{720(1-a)^3},
 \qquad \Lambda=Y+B.
\]
Then \eqref{eq:singular-master} gives the base inverse
\begin{equation}\label{eq:qone-infinite-t-inverse}
 t=\frac{A}{\Lambda}
 +\frac{A^2C}{\Lambda^3}
 +\frac{2A^3C^2-A^4E}{\Lambda^5}
 +\OO(\Lambda^{-7}),
 \qquad q=e^{-t}.
\end{equation}
In particular,
\begin{align}
 q={}&1-\frac{A}{\Lambda}+\frac{A^2}{2\Lambda^2}
 -\frac{A^3/6+A^2C}{\Lambda^3}+\OO(\Lambda^{-4}).
 \label{eq:qone-infinite-q-inverse}
\end{align}
The numerical table in the appendix shows the expected gain of two inverse
powers of $\Lambda$ at each correction.

\begin{remark}[Complex sectors and exponentially small terms]
Formula \eqref{eq:mcintosh} is a Poincar\'e expansion.  A complete global
inverse requires a branch of $\Polylogarithm_2(a)$ and $\PrincipalLogarithm(1-a)$ and, near Stokes
rays, exponentially small modular or resurgence corrections.  These effects
are invisible to algebraic powers of $1/\Lambda$ but can move the inverse by an
amount comparable with the least term of the expansion.
\end{remark}

\section{Fixed level, $a$-inversion as $q\to1$}
The previous regime holds $a$ fixed and sends $y$ to zero.  A different and
important inverse problem fixes $0<y<1$ and lets $a=a(t)$ tend to zero.  Put
$L=-\log y$.  Substituting a power series into the Lambert sum gives
\begin{align}
 a(t)={}&Lt-\frac{L(L+2)}4t^2
 +\frac{L(L^2+9L+12)}{72}t^3\notag\\
 &-\frac{L(L^3+4L^2+12L+24)}{576}t^4
 +\OO(t^5).
 \label{eq:a-fixed-level-qone}
\end{align}
This is a regular series in $t$, not in $1/\log y$.  It describes the moving
level curves that terminate at the corner $(a,q)=(0,1)$.

\begin{proof}[Derivation]
From \eqref{eq:lambert-log},
\[
 L=\sum_{m\ge1}\frac{a^m}{m(1-e^{-mt})}.
\]
Use
$1/(1-e^{-x})=x^{-1}+\tfrac12+x/12-x^3/720+\cdots$, insert
$a=c_1t+c_2t^2+\cdots$, and solve successively.  The first four equations give
\eqref{eq:a-fixed-level-qone}.
\end{proof}

\section{Radial asymptotics at every root of unity}
The singular point $q=1$ is one member of a family.  Let $\zeta$ be primitive
of order $d$ and approach it radially by
\[
 q=\zeta e^{-t},\qquad t\downarrow0.
\]
Group factors by their residue class modulo $d$:
\begin{equation}\label{eq:root-unity-factorization-infinite}
 (a;\zeta e^{-t})_\infty
 =\prod_{r=0}^{d-1}
 (a\zeta^re^{-rt};e^{-dt})_\infty.
\end{equation}

\begin{theorem}[Root-of-unity dilogarithm law]\label{thm:root-unity-dilog}
Let $|a|<1$ and fix compatible logarithm branches.  As $t\to0$ in a sector,
\begin{align}
 \PrincipalLogarithm(a;\zeta e^{-t})_\infty
={}&-\frac{\Polylogarithm_2(a^d)}{d^2t}+C_{d,\zeta}(a)+\OO(t),
 \label{eq:root-unity-asymptotic}\\
 C_{d,\zeta}(a)
={}&\sum_{r=0}^{d-1}\left(\frac12-\frac rd\right)
       \PrincipalLogarithm(1-a\zeta^r).\label{eq:root-unity-constant}
\end{align}
\end{theorem}
\begin{proof}
Apply \eqref{eq:mcintosh} with base $e^{-dt}$ to every factor in
\eqref{eq:root-unity-factorization-infinite}.  The leading term is
\[
 -\frac1{dt}\sum_{r=0}^{d-1}\Polylogarithm_2(a\zeta^r).
\]
The root-of-unity filter gives
\[
 \sum_{r=0}^{d-1}\Polylogarithm_2(a\zeta^r)
 =\frac1d\Polylogarithm_2(a^d).
\]
The parameter $a\zeta^re^{-rt}$ contributes
$-(r/d)\PrincipalLogarithm(1-a\zeta^r)$ from the first-order expansion of its dilogarithm,
while McIntosh contributes $\tfrac12\PrincipalLogarithm(1-a\zeta^r)$.  Summing gives
\eqref{eq:root-unity-constant}.
\end{proof}

\begin{corollary}[Base inverse near a root of unity]\label{cor:q-root-unity-inverse}
Let $y=(a;q)_\infty$ tend to zero along a branch for which
$A_d=\Polylogarithm_2(a^d)/d^2\ne0$.  Put $Y=-\PrincipalLogarithm y$ and
$\Lambda=Y+C_{d,\zeta}(a)$.  Then
\[
 t=\frac{A_d}{\Lambda}+\OO(\Lambda^{-3}),
 \qquad q=\zeta\exp\left(-\frac{A_d}{\Lambda}
 +\OO(\Lambda^{-3})\right).
\]
Higher terms follow by carrying \eqref{eq:root-unity-asymptotic} to higher
order and applying \eqref{eq:singular-master}.
\end{corollary}

\begin{corollary}[Universal fixed-level Puiseux scale]\label{cor:a-root-unity-scale}
Fix $y=e^{-L}$ with $L\ne0$.  On branches tending to $a=0$,
\begin{equation}\label{eq:a-root-unity-leading}
 a^d\sim d^2Lt,
 \qquad
 a\sim\omega(d^2Lt)^{1/d},\quad \omega^d=1.
\end{equation}
Thus radial inversion in $a$ has exponent $1/d$ at a primitive $d$-th root of
unity.
\end{corollary}
\begin{proof}
For small $a$, $\Polylogarithm_2(a^d)=a^d+\OO(a^{2d})$.  Insert this into
\eqref{eq:root-unity-asymptotic} while keeping $L$ fixed.
\end{proof}

The exponent $1/d$ is the analytic counterpart of the exact finite identity
$(a;\zeta)_{md}=(1-a^d)^m$.

\section{The radial point $q=-1$ in detail}
Take $d=2$, $\zeta=-1$.  The exact factorization is
\[
 (a;-e^{-t})_\infty
 =(a;e^{-2t})_\infty(-ae^{-t};e^{-2t})_\infty.
\]
For $0<a<1$, expansion to the first correction gives
\begin{equation}\label{eq:minus-one-infinite-fixed-a}
 \log(a;-e^{-t})_\infty
 =-\frac{\Polylogarithm_2(a^2)}{4t}+\frac12\log(1-a)
 -\frac{a(a+3)}{12(1-a^2)}t+\OO(t^3).
\end{equation}
There is no $t^2$ term.  If
\[
 A_-={\Polylogarithm_2(a^2)\over4},\qquad
 B_-={1\over2}\log(1-a),\qquad
 C_-={a(a+3)\over12(1-a^2)},
\]
then
\begin{equation}\label{eq:minus-one-inverse-t}
 t=\frac{A_-}{Y+B_-}
 +\frac{A_-^2C_-}{(Y+B_-)^3}+\OO(Y^{-5}),
 \qquad q=-e^{-t}.
\end{equation}

For the complementary fixed-level inverse, $a=\OO(t^{1/2})$.  Solving the
Lambert series in powers of $s=t^{1/2}$ gives the more precise positive branch
\begin{align}
 a={}&2\sqrt{Lt}-t
 +\frac{\tfrac14-L-L^2}{\sqrt L}t^{3/2}
 +\frac{4L+3}{6}t^2+\OO(t^{5/2}).
 \label{eq:a-fixed-level-minus-one}
\end{align}
The negative branch and the complex branches are obtained by the appropriate
choice of the leading square root, followed by coefficient recursion.

\section{Why $q=\pm\infty$ is different for the infinite product}
The product $(a;q)_\infty$ does not converge when $|q|>1$.  Therefore a
statement about its inverse at $q=\pm\infty$ is meaningless until a specific
regularization or reciprocal-base continuation is chosen.  Finite products
have the canonical algebraic theory of \cref{eq:q-infty-finite-first}; complex
orders can be continued using reciprocal-base definitions; but two different
regularizations of the infinite product can differ by exponential factors and
produce different inverse asymptotics.  This report consequently treats
$q=\pm\infty$ canonically for finite products and for functions such as
$\Gamma_q$ that possess a standard reciprocal-base definition.

\chapter{Inversion in the complex order and truncation index}

For $|q|<1$, define the complex-order shifted factorial by
\begin{equation}\label{eq:complex-order}
 (a;q)_\alpha=\frac{(a;q)_\infty}{(aq^\alpha;q)_\infty},
\end{equation}
with a chosen branch of $q^\alpha=e^{\alpha\PrincipalLogarithm q}$.

\section{Global real order inverse}
\begin{theorem}[Monotone order branch]\label{thm:alpha-global}
Let $0<a,q<1$.  The map
\[
 \alpha\longmapsto(a;q)_\alpha
\]
is strictly decreasing on $[0,\infty)$ from $1$ to $(a;q)_\infty$.
Consequently every $y\in((a;q)_\infty,1]$ has a unique real inverse order.
\end{theorem}
\begin{proof}
Differentiating \eqref{eq:complex-order} gives
\begin{equation}\label{eq:alpha-logder}
 \partial_\alpha\PrincipalLogarithm(a;q)_\alpha
 =\PrincipalLogarithm q\sum_{j\ge0}\frac{aq^{\alpha+j}}{1-aq^{\alpha+j}}<0.
\end{equation}
The endpoints follow directly from the definition.
\end{proof}

\section{Series at order $\alpha=0$}
Put $u=\PrincipalLogarithm y$ and
\[
 c_m=(\PrincipalLogarithm q)^m\sum_{j\ge0}\Polylogarithm_{1-m}(aq^j),\qquad m\ge1.
\]
Then
\[
 \PrincipalLogarithm(a;q)_\alpha
 =\sum_{m\ge1}\frac{c_m}{m!}\alpha^m,
\]
and hence
\begin{align}
 \alpha={}&\frac{u}{c_1}-\frac{c_2}{2c_1^3}u^2
 +\frac{3c_2^2-c_1c_3}{6c_1^5}u^3+\OO(u^4).
 \label{eq:alpha-near-zero-inverse}
\end{align}
The Lambert sums converge rapidly for fixed $q<1$ and give a practical local
inverse without numerical differencing.

\section{Large order and inversion of the tail}
Let
\[
 P_\infty=(a;q)_\infty,
 \qquad w=\frac{y}{P_\infty}-1,
 \qquad z=aq^\alpha.
\]
As $\alpha\to\infty$, Euler's reciprocal expansion gives
\[
 1+w=\frac1{(z;q)_\infty}
 =1+\frac{z}{1-q}
 +\frac{z^2}{(1-q)(1-q^2)}+\cdots.
\]
Reversion yields
\begin{align}
 z={}&(1-q)w-\frac{1-q}{1+q}w^2\notag\\
 &+\frac{(1-q)(2q^2+q+1)}{(1+q)^2(1+q+q^2)}w^3
 +\OO(w^4).
 \label{eq:z-tail-inverse}
\end{align}
Therefore
\begin{align}
 \alpha={1\over\PrincipalLogarithm q}\biggl[&\PrincipalLogarithm\frac{(1-q)w}{a}
 -\frac{w}{1+q}\notag\\
 &+\frac{3q^2+q+1}{2(1+q)^2(1+q+q^2)}w^2
 +\OO(w^3)\biggr].
 \label{eq:alpha-large-inverse}
\end{align}
This logarithmic divergence quantifies how many factors are needed before a
finite product is within a prescribed relative distance of the infinite one.

\section{Recovering a discrete truncation index}
For integer $n$ the sequence $(a;q)_n$ decreases to $P_\infty$.  Define
\[
 n_-(y)=\max\{n:(a;q)_n\ge y\},
 \qquad P_\infty<y<1.
\]
The continuous inverse \eqref{eq:alpha-large-inverse} gives an asymptotically
sharp estimate, and the exact discrete index is one of the two neighboring
integers.  A certified algorithm evaluates those two finite products, avoiding
any ambiguity caused by rounding a slowly varying logarithm.

\section{Joint inversion and level-curve differential equations}
For $F(a,q,\alpha)=(a;q)_\alpha$, a level surface $F=y$ satisfies
\[
 F_a\,da+F_q\,dq+F_\alpha\,d\alpha=0.
\]
Thus, for example,
\[
 \left.\frac{d\alpha}{dq}\right|_{a,y}
 =-\frac{F_q}{F_\alpha},
 \qquad
 \left.\frac{da}{dq}\right|_{\alpha,y}
 =-\frac{F_q}{F_a}.
\]
The logarithmic derivative representations make these ratios explicit Lambert
series.  They are useful both analytically, for proving monotonicity, and
numerically, as tangent predictors in continuation algorithms.

\part{Gaussian coefficients and $q$-gamma inverses}

\chapter{The base inverse of a Gaussian polynomial}

Let
\[
 B_{n,k}(q)=\GaussianBinomial{n}{k}{q},
 \qquad 0<k<n,
 \qquad D=k(n-k).
\]
It is a monic reciprocal polynomial of degree $D$ with nonnegative integer
coefficients:
\begin{equation}\label{eq:qbinom-reciprocity}
 B_{n,k}(q)=q^D B_{n,k}(q^{-1}).
\end{equation}
Its coefficient of $q^m$ counts partitions of $m$ fitting inside a
$k\times(n-k)$ rectangle.

\section{The unique positive inverse}
\begin{theorem}[Global positive Gaussian inverse]\label{thm:qbinom-positive-inverse}
The map $q\mapsto B_{n,k}(q)$ is strictly increasing from $1$ to $+\infty$ on
$[0,\infty)$.  Therefore every $y\ge1$ has a unique inverse
$Q_{n,k}(y)\ge0$.  For $y\ge\binom nk$,
\begin{equation}\label{eq:qbinom-global-bounds}
 \left(\frac{y}{\binom nk}\right)^{1/D}
 \le Q_{n,k}(y)\le y^{1/D}.
\end{equation}
\end{theorem}
\begin{proof}
Every nonconstant coefficient is nonnegative and the coefficient of $q$ is
$1$, so the derivative is positive for $q\ge0$.  For $q\ge1$, reciprocity and
positivity give
$q^D\le B_{n,k}(q)\le B_{n,k}(1)q^D=\binom nk q^D$.
Invert these inequalities.
\end{proof}

This theorem isolates a canonical sheet among the $D$ complex branches.  It is
also the best starting point for certified numerical inversion because
bisection never loses the branch.

\section{Partition-controlled reversion at $q=0$}
Write
\[
 B_{n,k}(q)=1+\sum_{m\ge1}p_{k,n-k}(m)q^m,
\]
where $p_{r,s}(m)$ counts partitions of $m$ in an $r\times s$ rectangle.  Put
$s=y-1$ and abbreviate
\[
 p_2=\mathbf1_{\{k\ge2\}}+\mathbf1_{\{n-k\ge2\}},
\]
\[
 p_3=\mathbf1_{\{n-k\ge3\}}
 +\mathbf1_{\{k\ge2,\ n-k\ge2\}}
 +\mathbf1_{\{k\ge3\}}.
\]
Then
\begin{equation}\label{eq:qbinom-qzero-inverse}
 q=s-p_2s^2+(2p_2^2-p_3)s^3+\OO(s^4).
\end{equation}
All higher coefficients are universal polynomials in the rectangle partition
numbers.  Lagrange--B\"urmann gives the compact all-orders formula
\begin{equation}\label{eq:qbinom-lagrange-qzero}
 [s^m]Q_{n,k}(1+s)
 =\frac1m[z^{m-1}]
 \left(\frac{z}{B_{n,k}(z)-1}\right)^m.
\end{equation}

\section{Bernoulli-cumulant inversion at $q=1$}
Set $q=e^t$, $C_0=\binom nk$, and
$u=\log(B_{n,k}(e^t)/C_0)$.  For
\[
 A_m(n,k)=\sum_{j=1}^k\bigl((n-k+j)^m-j^m\bigr),
\]
the exact formal expansion is
\begin{equation}\label{eq:qbinom-cumulants}
 u=\frac D2t+
 \sum_{r\ge1}\frac{B_{2r}A_{2r}(n,k)}{2r(2r)!}t^{2r}.
\end{equation}
In particular,
\[
 A_2=D(n+1),
 \qquad
 A_4=D(n+1)\bigl(n(n+1)-D\bigr).
\]
Reversion gives
\begin{align}
 t={}&\frac{2u}{D}-\frac{n+1}{3D^2}u^2
 +\frac{(n+1)^2}{9D^3}u^3\notag\\
 &+\frac{(n+1)\{6[n(n+1)-D]-25(n+1)^2\}}
 {540D^4}u^4+\OO(u^5),
 \label{eq:qbinom-qone-inverse}\\
 q={}&e^t.
\end{align}

\begin{proof}[Proof of the cumulant formula]
For every positive integer $m$,
\[
 \log\frac{e^{mt}-1}{m(e^t-1)}
 =\frac{m-1}{2}t+
 \sum_{r\ge1}\frac{B_{2r}(m^{2r}-1)}{2r(2r)!}t^{2r}.
\]
Apply this identity to the product formula for $B_{n,k}$ and sum over
$j=1,\ldots,k$.  The inverse series follows from ordinary reversion.
\end{proof}

The absence of odd cumulants beyond the linear term is a reciprocity effect.
Indeed \eqref{eq:qbinom-reciprocity} implies that
$e^{-Dt/2}B_{n,k}(e^t)$ is an even function of $t$.

\section{Puiseux inversion at $q=\pm\infty$}
Let $R=y^{1/D}$ on a chosen branch.  By reciprocity,
\[
 B_{n,k}(q)=q^D\left(1+q^{-1}+p_2q^{-2}+\OO(q^{-3})\right).
\]
Therefore
\begin{equation}\label{eq:qbinom-infty-inverse}
 q=R-\frac1D
 +\left(\frac{D-1}{2D^2}-\frac{p_2}{D}\right)R^{-1}
 +\OO(R^{-2}).
\end{equation}
The $D$ choices of $R$ are the $D$ inverse sheets at infinity.  A real
$q\to-\infty$ branch exists for those real levels whose sign agrees with
$q^D$; in particular, it reaches positive large $y$ when $D$ is even and
negative large $y$ when $D$ is odd.

The full expansion is algorithmic.  Substitute
$q=R+c_0+c_1R^{-1}+\cdots$ into
$q^DB_{n,k}(q^{-1})=R^D$ and solve recursively.  Only the first $m$
partition coefficients are required for the inverse through $R^{1-m}$.

\section{Exact local data at $q=-1$}
The point $q=-1$ is especially clean.

\begin{theorem}[Value and derivative at $-1$]\label{thm:qbinom-minus-one}
Write $n=2m$ or $2m+1$.
\begin{enumerate}
\item If $n=2m$ and $k=2r+1$, then
\begin{equation}\label{eq:qbinom-minusone-zero}
 B_{2m,2r+1}(-1)=0,
 \qquad
 B'_{2m,2r+1}(-1)
 =\frac{m!}{r!(m-r-1)!}.
\end{equation}
Thus the zero is simple and the inverse at $y=0$ is regular.
\item In every other nontrivial case,
\begin{equation}\label{eq:qbinom-minusone-value}
 B_{n,k}(-1)=\binom{\lfloor n/2\rfloor}{\lfloor k/2\rfloor},
 \qquad
 B'_{n,k}(-1)=-\frac D2B_{n,k}(-1).
\end{equation}
The inverse is again regular.
\end{enumerate}
\end{theorem}
\begin{proof}
The value is the $q$-Lucas theorem at a primitive square root of unity.  In the
nonzero case $D$ is even; differentiating reciprocity at $q=-1$ gives
$2B'(-1)=-DB(-1)$.  In the zero case, cancel the unique factor of
$1+q$ in the cyclotomic factorization and evaluate the remaining product,
which gives \eqref{eq:qbinom-minusone-zero}.
\end{proof}

Hence $q=-1$ itself creates no ramification for a nonconstant Gaussian
polynomial: every level occurring there has nonzero first derivative.  This
contrasts sharply with the Pochhammer family, where special $a$ values can
force quadratic or cubic criticality.

\section{Primitive roots of unity and cyclotomic valuation}
Let $\zeta$ be primitive of order $d$ and write
$n=ad+b$, $k=rd+s$, with $0\le b,s<d$.  The $q$-Lucas theorem states
\begin{equation}\label{eq:q-lucas}
 B_{n,k}(\zeta)=\binom ar B_{b,s}(\zeta).
\end{equation}
The exact exponent of the cyclotomic polynomial $\Phi_d(q)$ in $B_{n,k}$ is
\begin{equation}\label{eq:cyclotomic-exponent}
 e_d(n,k)=\left\lfloor\frac nd\right\rfloor
 -\left\lfloor\frac kd\right\rfloor
 -\left\lfloor\frac{n-k}{d}\right\rfloor\in\{0,1\}.
\end{equation}
Thus whenever $B_{n,k}(\zeta)=0$, the zero is simple.  The local inverse at
$y=0$ is therefore analytic.  At nonzero root-of-unity values, criticality is
not forced by cyclotomy; it is detected by the derivative of the residual
polynomial.

\section{The complex branch locus}
For fixed $n,k$, the finite critical values are the roots in $y$ of
\[
 \Disc_q(B_{n,k}(q)-y).
\]
The radius of convergence of a local inverse series is bounded by the nearest
critical value in the chosen $y$-plane.  Since the coefficients are real and
the polynomial is reciprocal, critical points occur in conjugate and
reciprocal patterns.  Determining their exact geometry is a finite algebraic
problem for each rectangle, but no general closed classification appears to
be known.  This motivates \cref{conj:qbinom-monodromy} below.

\chapter{Continuous Gaussian parameters}

For $0<q<1$, define
\begin{equation}\label{eq:continuous-qbinom}
 \mathcal B_q(\alpha,\beta)
 =\frac{\Gamma_q(\alpha+1)}
 {\Gamma_q(\beta+1)\Gamma_q(\alpha-\beta+1)}.
\end{equation}
This interpolates the Gaussian polynomial at integer parameters and makes
inversion in ``$n$'' and ``$k$'' an analytic problem.

\section{Inversion in the upper parameter}
Let $\psi_q=\partial_x\log\Gamma_q$.  Then
\begin{equation}\label{eq:continuous-alpha-derivative}
 \partial_\alpha\log\mathcal B_q
 =\psi_q(\alpha+1)-\psi_q(\alpha-\beta+1).
\end{equation}
Because $\psi_q$ is strictly increasing, this derivative is positive for
$\beta>0$ on the domain $\alpha>\beta-1$.

\begin{theorem}[Global continuous upper-parameter inverse]
Fix $0<q<1$ and $\beta>0$.  The function
$\alpha\mapsto\mathcal B_q(\alpha,\beta)$ is strictly increasing on
$(\beta-1,\infty)$ from $0$ to
\begin{equation}\label{eq:continuous-upper-limit}
 \frac{1}{(1-q)^\beta\Gamma_q(\beta+1)}.
\end{equation}
It therefore has a unique real inverse for every level in this open range.
\end{theorem}
\begin{proof}
Monotonicity follows from \eqref{eq:continuous-alpha-derivative}.  At the left
endpoint the factor $\Gamma_q(\alpha-\beta+1)$ has a pole.  At infinity, the
$q$-gamma quotient tends to $(1-q)^{-\beta}$.
\end{proof}

At a regular base point $(\alpha_0,y_0)$, put
$u=\log(y/y_0)$ and
\[
 c_m=\psi_q^{(m-1)}(\alpha_0+1)
      -\psi_q^{(m-1)}(\alpha_0-\beta+1).
\]
Then the inverse series is exactly \eqref{eq:universal-reversion} with
$x=\alpha-\alpha_0$.

\section{Inversion in the lower parameter and the central branch point}
Differentiation gives
\begin{equation}\label{eq:continuous-beta-derivative}
 \partial_\beta\log\mathcal B_q
 =\psi_q(\alpha-\beta+1)-\psi_q(\beta+1).
\end{equation}
The function is symmetric under $\beta\mapsto\alpha-\beta$, so it reaches its
unique maximum at $\beta=\alpha/2$.  Put
\[
 A=\frac\alpha2+1,
 \qquad x=\beta-\frac\alpha2,
 \qquad \mathcal B_0=\mathcal B_q(\alpha,\alpha/2).
\]
Taylor expansion gives
\begin{equation}\label{eq:beta-center-direct}
 -\log\frac{\mathcal B_q(\alpha,\alpha/2+x)}{\mathcal B_0}
 =\psi_q'(A)x^2+\frac{\psi_q'''(A)}{12}x^4+\OO(x^6).
\end{equation}
Hence, with $v=-\log(y/\mathcal B_0)$,
\begin{equation}\label{eq:beta-center-inverse}
 x=\pm\sqrt{\frac{v}{\psi_q'(A)}}
 \left(1-\frac{\psi_q'''(A)}{24\psi_q'(A)^2}v
 +\OO(v^2)\right).
\end{equation}
This is a structural, not accidental, square-root inverse: the two lower
parameters $\beta$ and $\alpha-\beta$ are indistinguishable.

\section{Continuous parameter inversion near $q=1$}
The $q$-gamma expansion in the next chapter has the form
\[
 \log\Gamma_{e^{-t}}(x)=\log\Gamma(x)+A_1(x)t+A_2(x)t^2+\OO(t^4).
\]
Therefore
\begin{align}
 \log\mathcal B_{e^{-t}}(\alpha,\beta)
={}&\log\binom{\alpha}{\beta}
 +\Delta A_1(\alpha,\beta)t
 +\Delta A_2(\alpha,\beta)t^2+\OO(t^4),
 \label{eq:continuous-qbinom-qone}
\end{align}
where
\[
 \Delta A_j=A_j(\alpha+1)-A_j(\beta+1)
              -A_j(\alpha-\beta+1).
\]
In particular,
\begin{equation}\label{eq:delta-A1}
 \Delta A_1=-\frac{\beta(\alpha-\beta)}2.
\end{equation}
If $\alpha_0$ solves the classical level equation
$\binom{\alpha_0}{\beta}=y$ and
\[
 g_\alpha=\psi(\alpha_0+1)-\psi(\alpha_0-\beta+1)\ne0,
\]
then the corresponding $q$-inverse is
\begin{equation}\label{eq:continuous-alpha-qcorrection}
 \alpha_q=\alpha_0+
 \frac{\beta(\alpha_0-\beta)}{2g_\alpha}t+\OO(t^2).
\end{equation}
Higher corrections follow mechanically from
\eqref{eq:implicit-second}.

\chapter{Inverse theory for Jackson's $q$-gamma function}

For $0<q<1$,
\begin{equation}\label{eq:qgamma-def}
 \Gamma_q(x)=(1-q)^{1-x}\frac{(q;q)_\infty}{(q^x;q)_\infty}.
\end{equation}
Its logarithmic derivative is
\begin{equation}\label{eq:qdigamma}
 \psi_q(x)=-\log(1-q)
 +(\log q)\sum_{j\ge0}\frac{q^{x+j}}{1-q^{x+j}}.
\end{equation}

\section{The two real $x$-inverse branches}
\begin{theorem}[Unique minimum of $\Gamma_q$]\label{thm:qgamma-minimum}
For every $0<q<1$, $\psi_q$ is strictly increasing from $-\infty$ to
$-\log(1-q)>0$ on $(0,\infty)$.  Hence there is a unique $x_q^*>0$ with
$\psi_q(x_q^*)=0$.  The function $\Gamma_q$ decreases on $(0,x_q^*)$,
increases on $(x_q^*,\infty)$, and tends to $+\infty$ at both endpoints.
Consequently:
\begin{itemize}
\item levels above $\Gamma_q(x_q^*)$ have exactly two positive real inverses;
\item the minimum level has one double inverse;
\item lower levels have no positive real inverse.
\end{itemize}
\end{theorem}
\begin{proof}
Differentiating \eqref{eq:qdigamma} gives
\[
 \psi_q'(x)=(\log q)^2
 \sum_{j\ge0}\frac{q^{x+j}}{(1-q^{x+j})^2}>0.
\]
The endpoint limits follow from the $j=0$ term and from dominated convergence.
The behavior of $\Gamma_q$ then follows by integration and from
\eqref{eq:qgamma-def}.
\end{proof}

Let $y_*=\Gamma_q(x_q^*)$.  Since $\psi_q'(x_q^*)>0$,
\[
 \log\frac{\Gamma_q(x)}{y_*}
 =\frac{\psi_q'(x_q^*)}{2}(x-x_q^*)^2+\OO((x-x_q^*)^3).
\]
Thus the two branches coalesce as
\begin{equation}\label{eq:qgamma-min-inverse}
 x-x_q^*=\pm
 \sqrt{\frac{2\log(y/y_*)}{\psi_q'(x_q^*)}}
 +\OO(\log(y/y_*)).
\end{equation}

\section{The limit $q\to1$ and inversion in $x$}
Put $q=e^{-t}$.  Moak's $q$-Stirling theory and the shifted-factorial
asymptotics imply \cite{Moak1984,McIntosh1999}
\begin{equation}\label{eq:qgamma-qone-expansion}
 \log\Gamma_{e^{-t}}(x)
 =\log\Gamma(x)+A_1(x)t+A_2(x)t^2+\OO(t^4),
\end{equation}
with no cubic term and
\begin{equation}\label{eq:qgamma-A12}
 A_1(x)=-\frac{(x-1)(x-2)}4,
 \qquad
 A_2(x)=\frac{B_3(x)}{72}-\frac{x-1}{24},
\end{equation}
where $B_3(x)=x^3-\tfrac32x^2+\tfrac12x$.

Fix a regular classical inverse $x_0$ satisfying $\Gamma(x_0)=y$ and
$\psi(x_0)\ne0$.  Write
$x_q=x_0+d_1t+d_2t^2+\cdots$.  Then
\begin{align}
 d_1={}&-\frac{A_1(x_0)}{\psi(x_0)},\label{eq:qgamma-xinverse-d1}\\
 d_2={}&-\frac{A_2(x_0)+A_1'(x_0)d_1
 +\tfrac12\psi'(x_0)d_1^2}{\psi(x_0)}.
 \label{eq:qgamma-xinverse-d2}
\end{align}
The companion experiment for $x_0=3.4$ shows second-order errors decreasing
as $t^3$.

\section{Inversion in the base near $q=1$}
Fix $x\notin\{1,2\}$ and put
$u=\log(y/\Gamma(x))$.  Because $A_1(x)\ne0$,
\begin{equation}\label{eq:qgamma-qinverse-qone}
 t=\frac{u}{A_1(x)}-\frac{A_2(x)}{A_1(x)^3}u^2
 +\OO(u^3),
 \qquad q=e^{-t}.
\end{equation}
At $x=1$ and $x=2$, $\Gamma_q(x)\equiv1$, so no base inverse exists.  These
are genuine constant degeneracies, not missing terms in
\eqref{eq:qgamma-qinverse-qone}.

\section{Inversion near $q=0$}
Expansion of \eqref{eq:qgamma-def} gives
\begin{equation}\label{eq:qgamma-qzero-direct}
 \log\Gamma_q(x)=q^x+(x-2)q
 +\OO\!\left(q^{\min(2,x+1,2x)}\right).
\end{equation}
The dominant term changes with $x$:
\begin{enumerate}
\item If $0<x<1$, then $q^x$ dominates and
\begin{equation}\label{eq:qgamma-qzero-inverse-subunit}
 q\sim(\log y)^{1/x}.
\end{equation}
This is a fractional-power inverse even when $x$ is irrational.
\item If $x>1$ and $x\ne2$, then
\begin{equation}\label{eq:qgamma-qzero-inverse-superunit}
 q\sim\frac{\log y}{x-2}.
\end{equation}
\item At $x=1,2$, the function is identically $1$.
\end{enumerate}
The full expansion lives in a Hahn series generated by the exponents $1$ and
$x$ when $x$ is nonintegral.

\section{The standard $q>1$ continuation and $q\to\infty$}
Jackson's reciprocal-base continuation satisfies
\begin{equation}\label{eq:qgamma-reciprocal-base}
 \Gamma_q(x)=q^{(x-1)(x-2)/2}\Gamma_{1/q}(x),
 \qquad q>1.
\end{equation}
Let
$A=(x-1)(x-2)/2$.  Since $\Gamma_{1/q}(x)\to1$ as $q\to\infty$,
\begin{equation}\label{eq:qgamma-infty-leading}
 \Gamma_q(x)\sim q^A.
\end{equation}
For $A\ne0$, the base inverse has the leading Puiseux form
\begin{equation}\label{eq:qgamma-infty-inverse}
 q\sim y^{1/A},
\end{equation}
with the branch selected by the reciprocal-base definition.  Corrections are
obtained by inserting the small-$1/q$ expansion
\eqref{eq:qgamma-qzero-direct} into \eqref{eq:qgamma-reciprocal-base}.  Again
$x=1,2$ are constant exceptions.

\part{Other $q$-analogs and transferable inverse schemes}

\chapter{$q$-numbers, $q$-factorials, and $q$-exponentials}

\section{The $q$-number}
For
\[
 [x]_q=\frac{1-q^x}{1-q},
\]
inversion in $x$ is elementary once logarithm branches are chosen:
\begin{equation}\label{eq:qnumber-xinverse}
 x=\frac{\PrincipalLogarithm(1-y(1-q))}{\PrincipalLogarithm q}.
\end{equation}
This exact formula is useful as a model: branch changes in $\PrincipalLogarithm$ create all
complex inverse sheets, while the positive real branch is unique for
$0<q<1$ and $y>0$ in its range.

For base inversion, integer $x=n$ gives the polynomial
$[n]_q=1+q+\cdots+q^{n-1}$.  Near $q=1$, set $q=e^t$ and
$u=\log([x]_{e^t}/x)$.  The Bernoulli expansion is
\begin{equation}\label{eq:qnumber-log-expansion}
 u=\frac{x-1}{2}t+
 \sum_{r\ge1}\frac{B_{2r}(x^{2r}-1)}{2r(2r)!}t^{2r}.
\end{equation}
It is inverted by \eqref{eq:universal-reversion}.  At infinity, for integer
$n\ge2$ and $R=y^{1/(n-1)}$,
\begin{equation}\label{eq:qnumber-infty-inverse}
 q=R-\frac1{n-1}+\OO(R^{-1}).
\end{equation}
For noninteger $x$, the equation
$q^x-yq+(y-1)=0$ defines a generalized algebraic-transcendental correspondence;
Newton polygons in the exponents $1$ and $x$ replace ordinary polynomial
analysis.

\section{The $q$-factorial}
Let
\[
 [n]_q!=\prod_{j=1}^n[j]_q,
 \qquad N=\frac{n(n-1)}2.
\]
At $q=0$,
\[
 [n]_q!=1+(n-1)q+\frac{(n-2)(n+1)}2q^2+\OO(q^3).
\]
Therefore, with $s=y-1$,
\begin{equation}\label{eq:qfactorial-qzero-inverse}
 q=\frac{s}{n-1}
 -\frac{(n-2)(n+1)}{2(n-1)^3}s^2+\OO(s^3).
\end{equation}
At $q=e^t$ near $1$,
\begin{equation}\label{eq:qfactorial-qone-direct}
 \log\frac{[n]_{e^t}!}{n!}
 =\frac N2t+
 \sum_{r\ge1}\frac{B_{2r}}{2r(2r)!}
 \sum_{j=1}^n(j^{2r}-1)t^{2r}.
\end{equation}
This is again a cumulant series with only one odd term.  Reciprocity gives
$[n]_q!=q^N[n]_{1/q}!$, and therefore
\begin{equation}\label{eq:qfactorial-infty-inverse}
 q=y^{1/N}-\frac{n-1}{N}+\OO(y^{-1/N}).
\end{equation}

\section{Euler's two $q$-exponentials}
Use
\[
 e_q(z)=\frac1{(z;q)_\infty},
 \qquad E_q(z)=(-z;q)_\infty.
\]
For $0<q<1$, $e_q$ maps $(-\infty,1)$ strictly increasingly onto
$(0,\infty)$, and $E_q$ maps $(-1,\infty)$ strictly increasingly onto
$(0,\infty)$.  Thus both possess global positive-real inverses.

Put $w=y-1$.  Euler's series and reversion give
\begin{align}
 e_q^{-1}(1+w)={}&(1-q)w-\frac{1-q}{1+q}w^2\notag\\
 &+\frac{(1-q)(2q^2+q+1)}{(1+q)^2(1+q+q^2)}w^3
 +\OO(w^4),\label{eq:eq-inverse}\\
 E_q^{-1}(1+w)={}&(1-q)w-\frac{q(1-q)}{1+q}w^2\notag\\
 &+\frac{q^2(1-q)(q^2+q+2)}{(1+q)^2(1+q+q^2)}w^3
 +\OO(w^4).\label{eq:Eq-inverse}
\end{align}
Under the classical scaling $z=(1-q)x$, both direct functions tend to $e^x$,
so their scaled inverses tend to $\log y$.

\section{A normalized $q$-Lambert inverse}
To obtain a deformation of the classical Lambert $W$ with a regular
$q\to1$ limit, define $\mathcal W_q$ locally by
\begin{equation}\label{eq:qLambert-def}
 y=x e_q((1-q)x)
 \quad\Longleftrightarrow\quad x=\mathcal W_q(y).
\end{equation}
Since
\[
 x e_q((1-q)x)=x+x^2+\frac{x^3}{1+q}+\OO(x^4),
\]
we obtain
\begin{equation}\label{eq:qLambert-series}
 \mathcal W_q(y)=y-y^2+\frac{1+2q}{1+q}y^3+\OO(y^4).
\end{equation}
As $q\to1^-$, the coefficient tends to $3/2$ and
$\mathcal W_q\to W$ locally.  The critical points of the direct map satisfy
\[
 1+(1-q)x\,\partial_z\log e_q(z)\big|_{z=(1-q)x}=0,
\]
and their images are the branch points of $\mathcal W_q$.  Mapping their
motion from $q=0$ to $q=1$ is a concrete continuation problem.

\chapter{Rogers--Szeg\H{o}, $q$-Catalan, theta, and basic hypergeometric inverses}

\section{Rogers--Szeg\H{o} polynomials}
Let
\[
 H_n(z;q)=\sum_{k=0}^n\GaussianBinomial{n}{k}{q}z^k.
\]
Inversion in $z$ at the unit level is regular when $[n]_q\ne0$:
\begin{equation}\label{eq:rogers-z-inverse}
 z=\frac{y-1}{[n]_q}
 -\frac{\GaussianBinomial{n}{2}{q}}{[n]_q^3}(y-1)^2
 +\OO((y-1)^3).
\end{equation}
For inversion in $q$, the direct endpoint data are
\begin{align}
 H_n(z;0)&=1+z+\cdots+z^n,\notag\\
 \partial_qH_n(z;0)&=\sum_{k=1}^{n-1}z^k,\label{eq:rogers-qzero-data}\\
 H_n(z;1)&=(1+z)^n,\notag\\
 \partial_qH_n(z;1)&=\frac{n(n-1)}2z(1+z)^{n-2}.
 \label{eq:rogers-qone-data}
\end{align}
Thus the base inverse is locally linear at either endpoint unless the displayed
slope vanishes.  When it does vanish, the next nonzero $q$-derivative determines
a Puiseux exponent.  At $q=\infty$, the Newton polygon is controlled by the
maximum of $k(n-k)$ among those $k$ for which $z^k\ne0$; a tie for odd $n$
can change the leading coefficient and create additional exceptional levels.

\section{MacMahon's $q$-Catalan polynomial}
Define
\[
 \Cat_n(q)=\frac{1}{[n+1]_q}\GaussianBinomial{2n}{n}{q}.
\]
For every $n\ge2$,
\begin{equation}\label{eq:qCatalan-qzero}
 \Cat_n(q)=1+q^2+\OO(q^3).
\end{equation}
The linear term cancels, so the inverse at $y=1$ is necessarily ramified:
\begin{equation}\label{eq:qCatalan-inverse}
 q=\pm\sqrt{y-1}+\OO(y-1).
\end{equation}
The positive sign is the canonical inverse on $q\ge0$.  At $q=1$,
\[
 \Cat_n(1)=\frac1{n+1}\binom{2n}{n},
 \qquad
 \frac{\Cat_n'(1)}{\Cat_n(1)}=\frac{n(n-1)}2,
\]
so the inverse is regular there for $n\ge2$.  This example shows that a ratio
of two individually regular $q$-analogs can acquire a critical endpoint by
coefficient cancellation.

\section{Theta products}
Let
\begin{equation}\label{eq:theta-def}
 \Theta(z;q)=(z;q)_\infty(q/z;q)_\infty,
 \qquad 0<|q|<1.
\end{equation}
Its zeros $z=q^m$, $m\in\IntegerNumbers$, are simple when $q$ is not a root of unity, so
the inverse in $z$ is analytic at the zero level.  The functional equations
\[
 \Theta(qz;q)=-z^{-1}\Theta(z;q),
 \qquad
 \Theta(q/z;q)=\Theta(z;q)
\]
organize the sheets.  The involution $z\mapsto q/z$ has fixed points
$z=\pm\sqrt q$; differentiation of the symmetry shows
\[
 \Theta_z(\pm\sqrt q;q)=0.
\]
Whenever the second derivative is nonzero, the inverse in $z$ has a
square-root branch at the corresponding theta value.  This is the theta
analogue of the central lower-parameter branch of the continuous Gaussian
coefficient.

The logarithmic derivative
\[
 \partial_z\log\Theta(z;q)
 =-\sum_{j\ge0}\frac{q^j}{1-zq^j}
 +\frac1z\sum_{j\ge0}\frac{q^{j+1}/z}{1-q^{j+1}/z}
\]
provides both local coefficients and a stable Newton iteration away from
zeros and critical points.

\section{A universal inverse template for ${}_r\phi_s$}
Write a basic hypergeometric function in a fixed convention as
\[
 F(z)={}_r\phi_s(\bm a;\bm b;q,z)
 =1+c_1z+c_2z^2+c_3z^3+\cdots,
 \qquad c_1\ne0.
\]
Then its local argument inverse is
\begin{equation}\label{eq:basic-hypergeom-inverse}
 z=\frac{y-1}{c_1}-\frac{c_2}{c_1^3}(y-1)^2
 +\left(\frac{2c_2^2}{c_1^5}-\frac{c_3}{c_1^4}\right)(y-1)^3
 +\OO((y-1)^4).
\end{equation}
Because each $c_m$ is a quotient of finite $q$-Pochhammer symbols, every
parameter derivative can be written in finite Lambert sums.  The same
procedure therefore inverts in any numerator or denominator parameter at a
regular point.  At terminating parameters, coalescing zeros can force
Puiseux series, and at $q\to1$ the correct scaling should be taken before
reversion to recover the corresponding ordinary hypergeometric inverse.

\section{Products and ratios with many parameters}
A broad class is
\begin{equation}\label{eq:multi-product}
 F(\bm a,q)=\prod_{\nu=1}^r(a_\nu;q)_{n_\nu}^{\eps_\nu},
 \qquad \eps_\nu\in\IntegerNumbers.
\end{equation}
For a selected parameter $a_\mu$,
\[
 \partial_{a_\mu}\log F
 =-\eps_\mu\sum_{j=0}^{n_\mu-1}
 \frac{q^j}{1-a_\mu q^j}.
\]
Thus all regular inverse coefficients follow from the universal calculus.
The singular set is the union of factor zeros, poles, and the critical
hypersurface on which the signed Lambert sums cancel.  This description
extends immediately to ratios defining orthogonal polynomials, $q$-beta
functions, and many terminating basic hypergeometric summands.

\chapter{Series algorithms, continuation, and certified computation}

\section{A branch-aware workflow}
A reliable inverse computation should separate symbolic branch analysis from
numerical evaluation.

\begin{algorithm}[Local and global inversion workflow]\label{alg:workflow}
Given $F(x;\lambda)=y$:
\begin{enumerate}
\item Choose a base point $(x_0,y_0,\lambda_0)$ and compute the first nonzero
      derivative of $F-y_0$ in $x$.
\item If the order is one, compute a Taylor inverse by Lagrange reversion or
      Newton iteration in a truncated series ring.  If the order is $m>1$,
      introduce $w=(y-y_0)^{1/m}$ and compute a Puiseux inverse.
\item Factor exact cyclotomic and Pochhammer zeros before numerical evaluation.
      Never evaluate a removable $0/0$ Gaussian product directly at a root of
      unity.
\item Near $q=1$ or another unit-circle root, replace $q$ by
      $\zeta e^{-t}$ and invert the logarithmic asymptotic in $t$.
\item Continue the selected branch with a tangent predictor from
      \eqref{eq:implicit-first} and a Newton corrector.
\item Monitor $|F_x|$.  When it becomes small, switch to a local Puiseux
      coordinate or use pseudo-arclength continuation.
\item Certify the result by an interval derivative bound, interval Newton, or
      a complex Rouch\'e test.
\end{enumerate}
\end{algorithm}

\section{Formal reversion without expression swell}
Direct substitution into a long symbolic product is expensive.  Better
methods are:
\begin{itemize}
\item work with logarithmic cumulants and Bell polynomials;
\item use Newton iteration in the truncated ring
      $\ComplexNumbers[[u]]/(u^{M+1})$, doubling the number of correct coefficients at each
      step;
\item exploit coefficient stability: a finite product with $n>M$ has the same
      $q=0$ inverse jet through order $M$ as the infinite product;
\item for Gaussian polynomials, obtain only the required rectangle partition
      numbers rather than expand the full polynomial;
\item at infinity, work in $s=1/q$ and use reciprocal symmetry.
\end{itemize}

If $g_0(u)$ is an initial inverse approximation, the truncated Newton update
\begin{equation}\label{eq:series-newton}
 g_{r+1}(u)=g_r(u)-\frac{F(g_r(u))-u}{F'(g_r(u))}
 \pmod{u^{2^{r+1}}}
\end{equation}
provides a fast general implementation.

\section{Stable direct evaluation}
Near $q=1$, multiplying factors is badly conditioned.  For $|a|<1$, use the
Lambert series \eqref{eq:lambert-log}; for finite products use compensated sums
of logarithms.  Near a root of unity, first group factors as in
\eqref{eq:root-unity-factorization-infinite}.  For Gaussian coefficients at a
root of unity, use cyclotomic factorization or the polynomial recurrence
instead of the quotient of two vanishing Pochhammer products.

\section{Numerical branch continuation}
Suppose $F(q)=y$ and a branch is known at $y_0$.  The differential equation
\[
 \frac{dq}{dy}=\frac1{F'(q)}
\]
provides a predictor.  Close to a simple critical point, use
$w=\sqrt{y-y_c}$ as the continuation parameter.  For polynomial inverses, all
branches can be tracked simultaneously by homotopy from a large level, where
the Puiseux roots are explicitly separated.  Matching at a moderate level
then identifies the monodromy permutation induced by loops around critical
values.

\section{What the companion program checks}
The distributed script \texttt{inverse\_q\_analogs\_experiments.py} uses
80-digit arithmetic and exact SymPy polynomials.  It checks:
\begin{itemize}
\item the finite/infinite stable $q=0$ reversion and every low-$n$ exception;
\item singular inversion of $(a;q)_\infty$ at $q\to1$;
\item radial inversion at $q\to-1$;
\item the fixed-level $a$-inverse at $q\to1$;
\item the Gaussian $q=1$ inverse and exact $q=-1$ values and derivatives;
\item the $q$-gamma inverse correction near $q=1$;
\item square-root ramification of the MacMahon $q$-Catalan polynomial;
\item the exact $x$-inverse of the $q$-number.
\end{itemize}
The output is reproduced in \cref{app:numerical-results}.

\part{Frontier questions and synthesis}

\chapter{Conjectures and research programs}

\section{Monodromy and generic Galois groups}
Polynomial inversion naturally asks for the permutation group generated by
continuation around the critical values.

\begin{conjecture}[Generic Pochhammer monodromy]\label{conj:poch-monodromy}
Let $n\ge3$ and $N=n(n-1)/2$.  Over the rational function field
$\ComplexNumbers(a,y)$, the polynomial
\[
 (a;q)_n-y
\]
viewed as a polynomial in $q$ has full symmetric geometric monodromy
$S_N$, except on an explicitly classifiable finite union of degenerate
parameter strata and low-degree cases.
\end{conjecture}

Evidence should be sought by computing discriminants, branch cycles, and
specializations with certified Galois groups.  The large-$y$ Puiseux branches
from \eqref{eq:s-lagrange-infty} give a natural base fiber for monodromy
tracking.

\begin{conjecture}[Generic Gaussian monodromy]\label{conj:qbinom-monodromy}
For rectangles outside an explicitly classifiable decomposable family, the
polynomial map
\[
 q\longmapsto\GaussianBinomial{n}{k}{q}
\]
has full symmetric monodromy $S_D$, $D=k(n-k)$, over a generic level.
Reciprocity constrains the critical configuration but does not force an
imprimitive inverse action for generic $y$.
\end{conjecture}

This formulation intentionally allows exceptions.  Composition identities,
small rectangles, and accidental critical-value collisions must be removed
before a full-symmetric statement can be true.

\section{All-orders root-of-unity inversion}
\begin{conjecture}[Cyclotomic polylogarithm closure]\label{conj:cyclotomic-closure}
For primitive $\zeta$ of order $d$, the complete radial expansion of
$\PrincipalLogarithm(a;\zeta e^{-t})_\infty$ has coefficients in the differential algebra
generated by
\[
 \Polylogarithm_m(a\zeta^r),\quad \PrincipalLogarithm(1-a\zeta^r),
 \qquad 0\le r<d,
\]
with rational functions of $a$ and cyclotomic constants.  After inversion, the
base series belongs to the corresponding algebra in odd inverse powers of the
large logarithmic coordinate, up to exponentially small sectors.
\end{conjecture}

The first part should follow systematically from Euler--Maclaurin expansion of
each factor in \eqref{eq:root-unity-factorization-infinite}.  The difficult
part is a uniform statement across branch cuts and Stokes directions.

\begin{problem}[Uniform $a\asymp t^{1/d}$ expansion]
Construct a two-parameter uniform expansion for
$(a;\zeta e^{-t})_\infty$ when $a=t^{1/d}A$ and $A$ ranges over compact
complex sets.  Determine the complete inverse series $A=A(y,t)$ and its
Stokes transitions.  The formulas \eqref{eq:a-fixed-level-qone} and
\eqref{eq:a-fixed-level-minus-one} are the first two cases.
\end{problem}

\section{Resurgence and exponentially small inverse corrections}
The McIntosh series is generally divergent at sufficiently high order.  Its
least term predicts an exponentially small error associated with modular or
saddle contributions.  Inverting a divergent expansion is not accomplished
by formal reversion alone: exponentially small changes in $\log(a;q)_\infty$
can be amplified by the inverse condition number.

\begin{problem}[Resurgent inverse $q$-Pochhammer]
Determine the Borel singularities and Stokes constants of the $q\to1$ and
$q\to\zeta$ expansions, then derive the induced transseries for the inverse
base $q(y)$.  Quantify how the inverse Stokes jump depends on
$\Polylogarithm_2(a^d)$ and on proximity to a critical level.
\end{problem}

A promising method is to combine modular transformations of eta/theta
specializations with exact WKB-style inversion of the logarithmic coordinate.

\section{Critical values, zero attractors, and inverse conditioning}
For fixed $(n,k)$, the critical points of a Gaussian polynomial and their
images can be computed exactly, but their large-rectangle distribution is
largely unexplored.

\begin{problem}[Limit distribution of Gaussian critical values]
Let $k/n\to\rho\in(0,1)$.  After the natural logarithmic scaling in $q$ and
in the level, determine the limiting distribution of critical points and
critical values of $\GaussianBinomial{n}{k}{q}$.  Relate it to the limit shape of
partitions in a rectangle and to saddle points of the associated free energy.
\end{problem}

Such a result would determine asymptotic radii of convergence for inverse
series at $q=0$ and $q=1$ and would identify the best domains for rational or
conformal acceleration.

\section{Coefficient structure of inverse series}
The first inverse coefficients often display positivity or alternating signs,
but exact computations show that naive global sign conjectures fail even for
small rectangles.  The right question is more refined.

\begin{problem}[Combinatorics of reversion coefficients]
Find signed combinatorial models for
\[
 [s^m]Q_{n,k}(1+s)
 \quad\text{and}\quad
 [x^m]Q_n((1-a)+(a^2-a)x;a)
\]
in terms of decorated partitions, plane trees, or cumulant assemblies.
Classify the first sign changes and determine whether natural rescalings yield
log-concavity, total positivity, or finite-order sign regularity.
\end{problem}

Lagrange inversion already expresses the coefficients as weighted tree sums;
the missing step is a weight-preserving interpretation adapted to rectangle
partitions or Pochhammer subset sums.

\section{Multivariate and elliptic extensions}
A basic hypergeometric identity supplies an algebraic or analytic hypersurface
in all of its parameters.  Multivariate inversion can therefore reveal
relations invisible in one-variable specializations.

\begin{problem}[Inverse basic hypergeometric varieties]
Develop a branch atlas for simultaneous inversion in $(z,q,a_1,\ldots)$ of
terminating ${}_r\phi_s$ functions.  Identify discriminant components caused
by termination, balancing, and root-of-unity resonance.  Extend the method to
elliptic shifted factorials and elliptic gamma functions, where two nomes
produce interacting unit-circle singularities.
\end{problem}

The proper local tool is the multivariate Weierstrass preparation theorem;
the proper global tool is numerical algebraic geometry combined with exact
special-function identities.

\section{Connections to geometric-comb interpolation and atomic functions}
The source monograph emphasizes Newton and Lagrange interpolation on geometric
nodes.  The basis factors of such interpolation are shifted products
\[
 \prod_{j=0}^{m-1}(x-q^j),
\]
which are elementary transforms of $q$-Pochhammer symbols.  Their inverse
level sets determine where the basis is well conditioned and where several
geometric nodes create ramification.  The complex-order inverse
\eqref{eq:alpha-large-inverse} also gives an ``effective depth'' of a geometric
tail: it answers how many dyadic or $q$-adic scales are needed to achieve a
given product level.

For the Rvachev up-function and the Fabius function, Fourier transforms and
functional equations involve products over dyadic scales.  Although the
factors are sinc functions rather than literal shifted factorials, the same
architecture applies: logarithmic cumulants for regular inversion, Lambert
$W$-type coordinates for exponentially flat limits, and Puiseux branches when
symmetry forces criticality.  A concrete research direction is to transfer the
certified branch-continuation machinery of this report to inverse level sets of
infinite sinc products and then compare their effective-depth variable with
$\alpha$ in \eqref{eq:complex-order}.

\section{Formal verification}
A realistic Lean development can be modular.
\begin{enumerate}
\item Formal power-series reversion and Lagrange--B\"urmann coefficients.
\item A local Puiseux theorem for a factored holomorphic germ, initially in a
      formal setting.
\item Finite $q$-Pochhammer derivative and zero formulas such as
      \eqref{eq:poch-a-derivative-root}.
\item Positivity and monotonicity proofs for the global real branches.
\item Cyclotomic factorization and the $q$-Lucas evaluation.
\item The exact $q=-1$ formulas in \cref{thm:critical-minus-one,thm:qbinom-minus-one}.
\item An asymptotic algebra supporting substitution and inversion of
      Poincar\'e series, culminating in \eqref{eq:qone-infinite-t-inverse}.
\end{enumerate}
The finite algebraic layers are suitable for immediate formalization.  The
uniform analytic and resurgent layers require a larger asymptotic library.

\chapter{Conclusions}
The inverse theory of $q$-analogs is not a collection of isolated numerical
root-finding problems.  It has a coherent architecture:
\begin{itemize}
\item products suggest logarithmic cumulants;
\item ordinary points give Taylor reversion;
\item zero multiplicities and symmetry give Puiseux exponents;
\item reciprocal polynomials give expansions at $\pm\infty$;
\item cyclotomic factorization controls roots of unity;
\item dilogarithmic singularities at the unit circle invert in a large
      logarithmic coordinate;
\item continuous parameters are governed by $q$-digamma monotonicity;
\item global real branches can often be isolated before any complex analysis.
\end{itemize}

The most novel synthesis is the root-of-unity law
\eqref{eq:root-unity-asymptotic}: it unifies $q\to1$ and radial $q\to-1$ and
predicts the fixed-level inverse exponent $1/d$ at every primitive $d$-th
root.  The finite theory mirrors it exactly through residue-block
factorization.  Together with the full branch atlas at infinity and the
continuous-parameter analysis, this gives a practical foundation for a much
larger program of inverse basic hypergeometric and geometric-scale functions.

\appendix

\chapter{Compact coefficient atlas}\label{app:atlas}

\section{Universal regular inverse}
If
$u=c_1x+c_2x^2/2!+c_3x^3/3!+c_4x^4/4!+\cdots$, then
\[
 x=\frac{u}{c_1}-\frac{c_2u^2}{2c_1^3}
 +\frac{(3c_2^2-c_1c_3)u^3}{6c_1^5}
 +\frac{(-c_1^2c_4+10c_1c_2c_3-15c_2^3)u^4}{24c_1^7}+\cdots.
\]

\section{Finite and infinite Pochhammer quick reference}
\begin{longtable}{>{\raggedright\arraybackslash}p{0.27\textwidth}
                  >{\raggedright\arraybackslash}p{0.65\textwidth}}
\toprule
Regime & Inverse expansion\\
\midrule
\endfirsthead
\toprule
Regime & Inverse expansion\\
\midrule
\endhead
$a\to0$, finite &
$a=z/E_1+E_2z^2/E_1^3+(2E_2^2/E_1^5-E_3/E_1^4)z^3+\cdots$, $z=1-y$\\
$q\to0$, finite/infinite &
$q=x-x^2+(a+1)x^3-(4a+1)x^4+(3a^2+11a+1)x^5+\cdots$\\
$q\to1$, finite &
$u=\sum_{m\ge1}-S_m(n)\Polylogarithm_{1-m}(a)t^m/m!$, followed by universal reversion\\
$q\to\infty$, finite &
$s=xH(s)^{1/N}$, $[x^m]s=m^{-1}[u^{m-1}]H(u)^{m/N}$, $q=1/s$\\
$a\to0$, infinite &
$a=(1-q)z+q(1-q)z^2/(1+q)+\cdots$\\
$q\to1$, infinite fixed $a$ &
$t=A/\Lambda+A^2C/\Lambda^3+(2A^3C^2-A^4E)/\Lambda^5+\cdots$\\
$q\to1$, infinite fixed $y=e^{-L}$ &
$a=Lt-L(L+2)t^2/4+L(L^2+9L+12)t^3/72+\cdots$\\
$q\to\zeta$, order $d$ &
$\log(a;q)_\infty=-\Polylogarithm_2(a^d)/(d^2t)+C_{d,\zeta}(a)+\OO(t)$\\
$a$ at fixed level near $\zeta$ &
$a\sim\omega(d^2Lt)^{1/d}$\\
complex order $\alpha\to\infty$ &
$\alpha=(\PrincipalLogarithm((1-q)w/a)-w/(1+q)+\cdots)/\PrincipalLogarithm q$\\
\bottomrule
\end{longtable}

\section{Gaussian quick reference}
\begin{longtable}{>{\raggedright\arraybackslash}p{0.25\textwidth}
                  >{\raggedright\arraybackslash}p{0.67\textwidth}}
\toprule
Regime & Inverse expansion\\
\midrule
\endfirsthead
\toprule
Regime & Inverse expansion\\
\midrule
\endhead
$q\to0$ & $q=s-p_2s^2+(2p_2^2-p_3)s^3+\cdots$, $s=y-1$\\
$q\to1$ & $t=2u/D-(n+1)u^2/(3D^2)+(n+1)^2u^3/(9D^3)+\cdots$\\
$q\to\infty$ & $q=y^{1/D}-D^{-1}+((D-1)/(2D^2)-p_2/D)y^{-1/D}+\cdots$\\
$q=-1$, zero case & $B'_{2m,2r+1}(-1)=m!/[r!(m-r-1)!]$\\
$q=-1$, nonzero case & $B'(-1)=-(D/2)B(-1)$\\
continuous $\beta$ at center &
$\beta-\alpha/2=\pm\sqrt{v/\psi_q'(\alpha/2+1)}(1+\OO(v))$\\
\bottomrule
\end{longtable}

\chapter{Numerical and symbolic validation}\label{app:numerical-results}
The following output is generated by the companion script.  The calculations
use 80 decimal digits; symbolic checks are exact.

\VerbatimInput[fontsize=\scriptsize]{numerical_results.txt}

\chapter{Companion-code notes}
The archive includes \texttt{inverse\_q\_analogs\_experiments.py}.  The script
is deliberately independent of the LaTeX source.  Its main implementation
choices are:
\begin{itemize}
\item the Lambert series
$-\sum_{m\ge1}a^m/[m(1-q^m)]$ for stable infinite-product logarithms;
\item exact Gaussian polynomials from the Pascal recurrence, avoiding
      removable root-of-unity singularities;
\item direct coefficient-by-coefficient symbolic reversion at $q=0$;
\item high-precision root solving only after an asymptotic formula supplies a
      branch-specific initial value;
\item assertions for every exact $q=-1$ value and derivative formula.
\end{itemize}
Run it from the report directory with
\begin{verbatim}
python inverse_q_analogs_experiments.py
\end{verbatim}
It overwrites \texttt{numerical\_results.txt} with a reproducible report.

\chapter{Bibliographic notes}
The direct $q$-product and Gaussian theory is classical; standard references
include Gasper--Rahman \cite{GasperRahman}, Andrews--Askey--Roy
\cite{AndrewsAskeyRoy}, Ismail \cite{Ismail}, and the NIST Digital Library of
Mathematical Functions \cite{DLMFq}.  McIntosh \cite{McIntosh1999} supplies
the shifted-factorial asymptotic central to the singular inverse at $q=1$;
Moak \cite{Moak1984} develops the $q$-Stirling framework.  Lagrange inversion,
analytic combinatorics, and branch singularities are treated in
\cite{Comtet,FlajoletSedgewick,Henrici}.  The repository sources
\cite{ProveItFrontiers,ProveItMonograph} provide the immediate context and
notation for this report.

\begin{thebibliography}{99}

\bibitem{ProveItFrontiers}
V.~Reshetnikov,
\emph{Exponents and $q$-Series Frontiers},
LaTeX source in the ProveIt repository,
\url{https://github.com/VladimirReshetnikov/ProveIt/tree/main/Analysis/FabiusFunction/docs/semi-formalized-research-frontiers/drafts/exponents-and-q-series/Exponents_and_q_Series_Frontiers},
accessed 30 August 2026.

\bibitem{ProveItMonograph}
V.~Reshetnikov,
\emph{$q$-Pochhammer Symbols and $q$-Binomial Coefficients: Algebra,
Combinatorics, Analysis, Arithmetic, and Summation Theory},
LaTeX source in the ProveIt repository,
\url{https://github.com/VladimirReshetnikov/ProveIt/tree/main/Analysis/FabiusFunction/docs/semi-formalized-research-frontiers/drafts/exponents-and-q-series/q_pochhammer_q_binomial_monograph},
accessed 30 August 2026.

\bibitem{GasperRahman}
G.~Gasper and M.~Rahman,
\emph{Basic Hypergeometric Series}, 2nd ed.,
Cambridge University Press, 2004.

\bibitem{AndrewsAskeyRoy}
G.~E.~Andrews, R.~Askey, and R.~Roy,
\emph{Special Functions},
Cambridge University Press, 1999.

\bibitem{Ismail}
M.~E.~H.~Ismail,
\emph{Classical and Quantum Orthogonal Polynomials in One Variable},
Cambridge University Press, 2005.

\bibitem{KacCheung}
V.~Kac and P.~Cheung,
\emph{Quantum Calculus},
Springer, 2002.

\bibitem{DLMFq}
NIST Digital Library of Mathematical Functions,
Chapter 17, ``$q$-Hypergeometric and Related Functions,''
\url{https://dlmf.nist.gov/17}, accessed 30 August 2026.

\bibitem{McIntosh1999}
R.~J.~McIntosh,
``Some asymptotic formulae for $q$-shifted factorials,''
\emph{Ramanujan Journal} \textbf{3} (1999), 205--214.
\url{https://doi.org/10.1023/A:1006949508631}.

\bibitem{Moak1984}
D.~S.~Moak,
``The $q$-analogue of Stirling's formula,''
\emph{Rocky Mountain Journal of Mathematics} \textbf{14} (1984), 403--414.
\url{https://doi.org/10.1216/RMJ-1984-14-2-403}.

\bibitem{Comtet}
L.~Comtet,
\emph{Advanced Combinatorics},
Reidel, 1974.

\bibitem{FlajoletSedgewick}
P.~Flajolet and R.~Sedgewick,
\emph{Analytic Combinatorics},
Cambridge University Press, 2009.

\bibitem{Henrici}
P.~Henrici,
\emph{Applied and Computational Complex Analysis}, Vol.~1,
Wiley, 1974.

\bibitem{StanleyEC1}
R.~P.~Stanley,
\emph{Enumerative Combinatorics}, Vol.~1, 2nd ed.,
Cambridge University Press, 2012.

\bibitem{CorlessLambertW}
R.~M.~Corless, G.~H.~Gonnet, D.~E.~G.~Hare, D.~J.~Jeffrey, and D.~E.~Knuth,
``On the Lambert $W$ function,''
\emph{Advances in Computational Mathematics} \textbf{5} (1996), 329--359.

\end{thebibliography}

\end{document}
